% user_guide.tex — complete reference for every user-facing feature
\chapter{User Guide}
\label{chap:user-guide}

This chapter is the authoritative reference for every user-facing command,
configuration option, document format, and error condition exposed by
\cgspkg{}.

% -----------------------------------------------------------------------
\section{CLI Overview}

The single entry-point is the \cgscli{} command.

\textbf{\cgscli{} is a Pixi task, not a globally installed executable.}
This project is developed and run exclusively through Pixi (\texttt{pixi
install} once, per checkout) --- there is no \texttt{pip install} that puts
a bare \texttt{cgitsync} on \texttt{\$PATH}. Every invocation shown in this
chapter as \texttt{pixi run cgitsync ...} must be typed exactly that way;
a bare \texttt{cgitsync ...} will not be found by the shell.

\begin{lstlisting}[language=bash]
$ pixi run cgitsync <command> [options]
\end{lstlisting}

Running \cgscli{} without a command prints a help summary and exits with
code~0.

\subsection{Prerequisites}

Three tools are needed before the first command: Git clones this checkout
and the repositories in a tree; Pixi installs the locked Python environment
and runs \cgscli{}; and repository creation needs the provider's own command
line tool --- \texttt{gh} for GitHub, \texttt{glab} for GitLab, or
\texttt{tea} for Codeberg. Install only the provider tool you use.

The complete list depends on the tree and the work being done:

\begin{center}
\begin{tabular}{lp{5.0cm}p{5.2cm}}
\textbf{Dependency} & \textbf{When it is needed} & \textbf{Where its version or requirement comes from} \\
\hline
Git & Every clone, inspection and synchronisation. & Its observed version is recorded with the State. \\
Pixi & Install this project's environment and run \cgscli{}. & \texttt{pixi.lock} pins the project environment; the observed Pixi version is recorded separately. \\
Python & Run \cgspkg{} inside the Pixi environment. & \texttt{pixi.toml} selects the interpreter; the observed interpreter is recorded. \\
\texttt{gh}, \texttt{glab}, or \texttt{tea} & \texttt{repo create}, according to the repository provider. & The provider tool reports its version and whether it is authenticated. \\
SSH client and agent & Repositories whose remote uses SSH. & Presence is observed. Keys, user names, and key paths are never recorded. \\
DVC or Git LFS & A tree that uses the corresponding data backend. & The tree declares its use; the observed tool version is recorded when used. \\
Compilers and system libraries & Native builds or software installed outside Pixi. & The tree declares these in its \texttt{.cgs} \texttt{[environment]} table. \\
\texttt{latexmk} and TeX & Developers rebuilding the manuals under \texttt{docs/}. & A developer requirement, not an end-user runtime dependency. \\
\end{tabular}
\end{center}

The machine must also match a platform supported by the tree's lock files.
\texttt{pixi run cgitsync env} reports the operating system, architecture,
Pixi platform, libc, tool versions, credential facts, and manifests it
observes. \texttt{pixi run cgitsync env check} compares that observation
with the tree's declared requirements. Neither command installs anything.

\cgspkg{} stores no credentials of its own and has no private-repository
authentication story: every git operation --- \texttt{clone}, \texttt{pull},
\texttt{push}, and the remote lookups behind \texttt{checkout}'s branch/tag
fallback --- runs against whatever the ambient environment already has
cached (\texttt{ssh-agent}, an HTTPS credential helper, a stored
personal-access-token). When that ambient credential is missing or stale,
the operation fails immediately with a normal, catchable error rather than
sitting silently on a terminal or GUI credential prompt that will never
arrive --- every git subprocess \cgspkg{} runs sets
\texttt{GIT\_TERMINAL\_PROMPT=0} and a no-op \texttt{GIT\_ASKPASS} for
exactly this reason. Fix the ambient credential (switch the remote to
\texttt{ssh}, refresh the token) and re-run the command; nothing needs to
be undone first.

% -----------------------------------------------------------------------
\section{Lifecycle Contract}

This guide follows the simplified canonical lifecycle:

\begin{enumerate}
  \item \texttt{initialise(.cgs)} $\rightarrow$ clone all repos $\rightarrow$
        \texttt{READY}  \textit{(new project)}\\
        \texttt{initialise(.gts)} $\rightarrow$ restore from snapshot $\rightarrow$
        \texttt{READY}  \textit{(existing project)}
  \item \texttt{pull(.cgs|.gts)} $\rightarrow$ resync an existing tree
  \item \texttt{checkout(.gts)} / \texttt{add} / \texttt{commit} / \texttt{push} / \texttt{tag}
        $\rightarrow$ git operations on the tree
  \item \texttt{freeze} $\rightarrow$ emit the next \texttt{.gts} id and update \texttt{<Project\_name>.lgr}
\end{enumerate}

The project-local \texttt{.lgr} register and sync ledger are both active (T24,
T32).  Every synchronisation operation is automatically recorded as an immutable
DAG event in the \texttt{[[ledger]]} section of the \texttt{.lgr} file.

% -----------------------------------------------------------------------
\section{Commands}

The CLI exposes three command groups once a \texttt{GitTree} is \texttt{READY}.
The minimalist level is \texttt{initialise}, \texttt{bootstrap},
\texttt{clean-init}, \texttt{freeze-release}, \texttt{freeze-release-force},
\texttt{status}, \texttt{view-tree}, and \texttt{launch-release}.  The expert level exposes
the individual operations: \texttt{branch}, \texttt{checkout}, \texttt{purge},
\texttt{validate}, \texttt{clone}, \texttt{pull}, \texttt{pull-force},
\texttt{autofix},
\texttt{add}, \texttt{commit}, \texttt{push}, \texttt{tag}, \texttt{freeze},
\texttt{import-submodules}, \texttt{init-from-submodules}, and
\texttt{verify}, plus \texttt{env} and \texttt{env check}.  The configuration level, \texttt{discover},
\texttt{configure} and
\texttt{create-cgs}, authors a \texttt{.cgs} specification without requiring
existing state.  Redundant
inspection aliases such as \texttt{print}, \texttt{view-operation}, and
\texttt{validate-topology} are not part of the public CLI surface.

\subsection{\texttt{initialise}}

\begin{lstlisting}[language=bash]
$ pixi run cgitsync initialise <spec.cgs> [--output-path DIR] [--commit-gitignore]
    [--force-gitignore-sync] [--git-user-name NAME] [--git-user-email EMAIL]
    [--force-protocol {ssh,https}] [--force-reclone]
$ pixi run cgitsync initialise <snapshot.gts>
$ pixi run cgitsync initialise --project NAME --repo PROVIDER:OWNER/REPO [--repo ...] [--output-path DIR]
\end{lstlisting}

Unified initialisation entry point.  Dispatches on file extension:

\begin{itemize}
  \item \texttt{.cgs} source (clone mode): clones the full repository tree.
    The optional \texttt{--output-path} controls
    \texttt{CGSPATH}; \texttt{CGSHOME} is derived as
    \texttt{CGSPATH/<project-name>}. When omitted, \texttt{CGSPATH} defaults to
    \texttt{../..}.  On success, ends in \texttt{READY}.  If clone setup fails
    because stale generated state is still present, the CLI prints
    \texttt{Try clean-init method}.
  \item \texttt{.gts} source (restore mode): loads a saved snapshot directly.
    On success, ends in the lifecycle state recorded in the snapshot.
  \item Direct CLI authoring: requires \texttt{--project} and at least one
    repeatable \texttt{--repo}. It is mutually exclusive with a source file.
    The CLI collects values only and calls
    \texttt{ComplexGitSyncClient.configure()}; that reusable Python API
    delegates parsing, normalization, and validation to \texttt{cgs\_format.py}
    and produces the same canonical \texttt{CgsDocument} as an equivalent
    \texttt{.cgs} file.
\end{itemize}

All initialization paths end in a \texttt{READY} tree or fail explicitly.

\paragraph{\texttt{initialise} re-clones your dependencies.} Only the root
repository at \texttt{CGSHOME} is kept as it is. Every repository below it
whose directory already holds files is \textbf{deleted and cloned again}. The
old \texttt{.git} goes with it, so nothing in that directory can be recovered
once the run has passed that point.

Before deleting anything, \texttt{initialise} inspects each destination and
stops the whole run when one of them holds work that exists nowhere else:
uncommitted changes, commits that have not been pushed to the branch's
upstream, or a branch with no upstream at all. It names every repository that
blocked it, in one message, and deletes none of them. The two ways forward
are to commit and push that work, or to pass \texttt{--force-reclone} and
delete it deliberately.

A destination that is not a Git checkout at all --- what a clone interrupted
halfway leaves behind --- is still cleared without any flag. That is the case
the deletion was written for, and it still works.

Before a \texttt{.cgs} source is confirmed ready, every repository with
children (root, or any nested repository that itself has further nested
children) is safely pulled parent-first and has its \texttt{.gitignore}
updated with the relative path of each immediate child --- nested
repositories are plain independent clones rather than gitlinks, so without
this a parent's \texttt{git status}/\texttt{git add} would otherwise see a
child's working tree as ordinary untracked content. This is logged as the
\texttt{GT-GITIGNORE} phase, after \texttt{GT-CLONE}. If the safe pull for
one of these repositories fails, \texttt{initialise} raises immediately and
nothing is written; pass \texttt{--force-gitignore-sync} to instead fall
back to a pull-force recovery for that one repository (never a force-push).
By default nothing is staged, committed, or pushed by this step --- it only
writes the file and prints what changed. Pass \texttt{--commit-gitignore}
to explicitly approve staging (\texttt{.gitignore} alone), committing (with
a message listing exactly which children were added), and pushing each
changed repository. The commit identity defaults to local git config;
\texttt{--git-user-name}/\texttt{--git-user-email} override it and persist
the override to \texttt{CGSHOME/.cgitsync/master.toml} via
\texttt{MasterConfig}, so later invocations on the same workspace pick it up
without repeating the flags.

The root's own \texttt{.gitignore} additionally always gains
\texttt{.cgitsync/} (the generated runtime-state directory --- snapshots,
register, run logs, never part of the \texttt{.cgs}/\texttt{.gts} project
spec) and its \texttt{<project-name>.lgr} register file, whether or not the
tree has any nested repositories at all --- that state is written under the
root regardless of tree shape, so it is excluded the same idempotent way,
via the same underlying write \texttt{sync\_gitignore()} already uses for
child mount points.

\texttt{--force-protocol \{ssh,https\}} is an expert option, typically used
systematically in CI rather than for everyday development. It overrides
\texttt{access\_protocol} in memory for every repository this run clones ---
including ones discovered later from a nested \texttt{.cgs} living in a
different, separately-cloned repository --- so no \texttt{.cgs} file is ever
read differently or written to. Forcing \texttt{https} is the usual CI case:
it removes the dependency on an SSH key or agent being present on the
runner, a dependency that otherwise differs between a \texttt{push} run and
a \texttt{pull\_request} run (GitHub Actions withholds repository secrets
from a pull request opened from a fork). Omitted (the default), each
repository follows whatever protocol its own \texttt{.cgs} entry declares,
exactly as without the flag.

\subsection{\texttt{bootstrap}}

\begin{lstlisting}[language=bash]
$ pixi run cgitsync bootstrap <spec.cgs> <project_name> [--cgs-path DIR]
    [--force-protocol {ssh,https}]
\end{lstlisting}

Clones the full repository tree from a \texttt{.cgs} source into an
isolated \texttt{CGSHOME}, for running \cgscli{} from its own standalone
clone rather than nested inside the project tree it manages --- one
\cgspkg{} clone reused across many projects. It delegates to the same
\texttt{ComplexGitSyncClient.clone\_cgs} pipeline as \texttt{clone} below,
and so shares its branch-fallback and nested-discovery behaviour, but
differs from both \texttt{clone} and \texttt{initialise} in two ways:

\begin{itemize}
  \item \texttt{project\_name} is a required positional argument, not
    inferred from the \texttt{.cgs} document. It always forms the final
    path segment of \texttt{CGSHOME}, regardless of the specification's own
    \texttt{project} name.
  \item \texttt{CGSPATH} does not default to \texttt{../..} relative to the
    current directory (that default assumes \cgspkg{} is cloned inside the
    tree it manages --- the nested-mode case \texttt{initialise} is for).
    When \texttt{--cgs-path} is omitted, \texttt{CGSPATH} instead defaults
    to a fresh \texttt{\$HOME/.cgs/CGS<timestamp>/} directory, created
    automatically if it does not exist, so a bootstrapped project can never
    land inside \cgspkg{}'s own clone.
\end{itemize}

\texttt{--force-protocol \{ssh,https\}} behaves exactly as described under
\texttt{initialise} above, applying here to every repository the run
clones --- including the root, which \texttt{bootstrap} (unlike
\texttt{initialise}) also clones from scratch rather than assuming already
checked out.

\subsection{\texttt{discover}}

\begin{lstlisting}[language=bash]
$ pixi run cgitsync discover [ROOT] [--write FILE] [--max-depth N]
\end{lstlisting}

Scans \texttt{ROOT} (default: the current directory) for git repositories
and drafts a \texttt{.cgs} from what is checked out, through the
\texttt{ComplexGitSyncClient.discover\_repos()} facade. This is the entry
point for adopting a project that exists on disk but has no \texttt{.cgs}
describing it yet.

The command is a pure read of the filesystem and of each repository's git
config: nothing is cloned, fetched, staged, or modified, and no network
call is made. Without \texttt{--write} it only prints what it found,
matching the ``report first, write only when asked'' posture of
\texttt{--commit-gitignore} and \texttt{import-submodules --apply}.

The walk descends with no depth limit by default, treating every directory
that contains a \texttt{.git} entry as a repository. Pass
\texttt{--max-depth N} (\texttt{ROOT} itself is depth~0) to bound it. It
never descends \emph{into} a \texttt{.git} directory: for a submodule the
real git directory lives at \texttt{<parent>/.git/modules/<name>} while the
child's own \texttt{.git} is a file, so walking into it would report the
same repository twice.

For each repository, \texttt{origin}'s URL is read and converted to the
canonical \texttt{provider:owner/repository} shorthand through the existing
\texttt{parse\_repo\_id()} grammar --- this command adds no second parser.
\texttt{relative\_path} is taken straight from the walk rather than from a
repository name, so a child mounted at \texttt{external/Thing} is recorded
there and not at \texttt{Thing}. \texttt{nested\_config} is left unset on
every entry: the default \texttt{auto} already resolves cleanly whether or
not a repository carries its own \texttt{*.cgs}.

Repositories with no \texttt{origin}, or whose remote URL does not resolve
to a registered provider, are reported as warnings and left out of the
draft. They are never guessed at: a draft that silently invented an address
would be worse than one that states what it could not determine. Always
review the generated file, then \texttt{validate} it.

A repository found \emph{inside} another repository is drafted as that
repository's child, not the project root's. Nesting is worked out by
comparing paths: every path in a \texttt{.cgs} is written from the project
root, so a repository whose path lies under another's sits inside it, and
the innermost one wins. The report marks such an entry
\texttt{inside: <path>} and prints the resulting tree, and
\texttt{build\_registry\_from\_cgs\_document()} reads the drafted file back
the same way --- the holding repository becomes a \texttt{PARENT} and keeps
its own \texttt{.gitignore} entry for the repository it holds, which is
what stops a later \texttt{add} from restaging that repository as a
submodule.

Only what is \emph{checked out at scan time} can be found. A repository
cloned without \texttt{--recurse-submodules} leaves its submodule paths as
empty directories, and those are correctly not reported; recovering them
from git metadata instead is \texttt{import-submodules}' job. The same
applies when \texttt{--max-depth} is given: if it stops the walk while
directories are still left to look into, the report says so in a warning
rather than presenting a partial answer as a complete one. Without the
flag, the walk cannot stop early --- there is nothing to stop it at.

\subsection{\texttt{configure}}

\begin{lstlisting}[language=bash]
$ pixi run cgitsync configure [--output FILE]
\end{lstlisting}

Interactively prompts for a project name, default branch, and one or more
repositories (provider, owner, name, and per-repository overrides), then
writes the resulting \texttt{.cgs} file through the same offline
\texttt{ComplexGitSyncClient.configure()} facade used by
\texttt{initialise --project/--repo} and \texttt{create-cgs}. The provider
prompt accepts \texttt{github}, \texttt{gitlab}, \texttt{codeberg}, or
\texttt{custom}; \texttt{custom} additionally prompts for the provider URL.
When \texttt{--output} is omitted, the CLI prompts for a destination path,
defaulting to \texttt{<project-name>.cgs}.

\subsection{\texttt{create-cgs}}

\begin{lstlisting}[language=bash]
$ pixi run cgitsync create-cgs --project CGSil1 \
    --repo github:flipoyo/ComplexGitSync \
    --repo codeberg:GX4G/GX4G \
    --output CGSil1.cgs
\end{lstlisting}

Builds the canonical \texttt{CgsDocument} through the offline format pipeline
and serializes concise \texttt{.cgs} TOML. It performs no Git or network work.

\subsection{\texttt{clean-init}}

\begin{lstlisting}[language=bash]
$ pixi run cgitsync clean-init <spec.cgs> [--output-path DIR] [--commit-gitignore]
    [--force-gitignore-sync] [--git-user-name NAME] [--git-user-email EMAIL]
    [--force-protocol {ssh,https}]
\end{lstlisting}

Recovery-oriented initialisation for a \texttt{.cgs} source. It follows the
same high-level pipeline as \texttt{initialise(.cgs)} (including the
\texttt{.gitignore} lifecycle sync and \texttt{--force-protocol} behaviour
described there), but inserts a cleanup phase before cloning:

\begin{center}
\texttt{load->expand->validate->purge->clone->gitignore}
\end{center}

Use it when an earlier initialisation attempt left partial child checkouts.
\texttt{CGSHOME/.cgitsync/master.toml} (the persisted Git identity override,
if any) is preserved by the purge phase --- it is workspace-local
configuration, not generated clone state.

\subsection{\texttt{purge}}

\begin{lstlisting}[language=bash]
$ pixi run cgitsync purge <spec.cgs> [--output-path DIR]
\end{lstlisting}

Removes generated clone state from the resolved \texttt{CGSHOME} without
cloning. It deletes child repository directories declared directly under the
project root and project-local \texttt{*.lgr} files. Nested directories that
are not immediate root children, and \texttt{.cgitsync/master.toml} (see
\texttt{clean-init} above), are left untouched by this command.

\subsection{\texttt{pull}}

\begin{lstlisting}[language=bash]
$ pixi run cgitsync pull [spec.cgs|snapshot.gts] [--commit-gitignore]
    [--force-gitignore-sync] [--git-user-name NAME] [--git-user-email EMAIL]
    [--force-protocol {ssh,https}]
\end{lstlisting}

Resynchronises the project tree.  \texttt{.cgs} sources reload the source,
discover nested configs, then resynchronise the loaded tree.  \texttt{.gts}
sources load the saved registry as the starting point and then run the same
pull workflow.  In both cases the operation starts at the project root:
\texttt{root -> parent -> leaf}.  Every repository --- root, parent, and leaf
alike --- is a plain independent clone and receives its own
\texttt{git pull --ff-only}.  Use \texttt{launch-release} when you want to
check out a frozen recorded release state instead of pulling current
remotes.  If local changes block this safe pull, the CLI prints
\texttt{You can try cgitsync pull-force command}.

Before pulling your own branch, \texttt{pull} fetches \emph{every} branch
of each remote. That is what lets \texttt{checkout} join a branch somebody
else pushed rather than starting a new one of the same name (see
\texttt{checkout} below). It also repairs a workspace whose clone left a
single-branch fetch refspec behind, described under \texttt{clone}.

For a \texttt{.cgs} source, \texttt{pull} also runs the same
\texttt{.gitignore} lifecycle sync described under \texttt{initialise} ---
including \texttt{--commit-gitignore}, \texttt{--force-gitignore-sync}, and
the \texttt{--git-user-name}/\texttt{--git-user-email} identity override ---
once the tree-wide pull completes. A \texttt{.gts} source runs no discovery,
so there is nothing new for the sync to find. \texttt{pull-force} does not
run this sync at all; it is a destructive recovery command, not a lifecycle
path the sync is wired into.

\texttt{--force-protocol \{ssh,https\}} rewrites every repo's remote to
that protocol before pulling, persisting the change (\texttt{git remote
set-url}, a no-op when the URL already matches) rather than a one-off
override --- the switch sticks for every command after this one too, the
same way a repo's protocol at clone time sticks for everything downstream
of it. Same meaning as \texttt{initialise}/\texttt{bootstrap}'s
\texttt{--force-protocol}, applied to an already-cloned tree instead of at
clone time --- useful when a repo was adopted rather than cloned by
\cgspkg{} (\texttt{initialise} never touches an adopted root's remote) and
its remote's protocol does not match what the ambient environment can
actually authenticate with. On a failure that looks like an SSH or HTTPS
auth problem, the error gains an actionable hint naming
\texttt{--force-protocol <the other one>} even without the flag.

The rewrite reads the repository's \emph{actual current} remote URL and
only swaps its scheme --- \texttt{ssh} (\texttt{user@host:path}) to/from
\texttt{https}/\texttt{http} (\texttt{scheme://host/path}) --- keeping the
host and path exactly as they already are. It does \emph{not} rebuild the
URL from the repository's stored identity fields
(\texttt{gitprovider}/\texttt{project\_owner\_name}/\dots{}): those can be
missing or stale for a repository loaded from an older \texttt{.gts}
snapshot (see the ``Formal \texttt{.gts} snapshot specification'' section
of the architecture reference), and rebuilding from a wrong or
default-GitHub provider would silently aim the pull at the wrong host.
Converting the URL that is already configured never has that failure mode.
A remote in neither recognised form (a bare local filesystem path, for
instance) cannot have its protocol converted; \texttt{--force-protocol}
then fails with an explicit error naming the repository, rather than
guessing.

\subsection{\texttt{pull-force}}

\begin{lstlisting}[language=bash]
$ pixi run cgitsync pull-force [spec.cgs|snapshot.gts] [--force-protocol {ssh,https}]
\end{lstlisting}

Destructively resynchronises the same tree.  Every repository --- root,
parent, and leaf alike --- is forced to the selected remote branch,
parent-first, with \texttt{git fetch}, \texttt{git checkout -B <branch>
FETCH\_HEAD}, and \texttt{git clean -fd}.  This command can discard local
uncommitted changes and untracked files; it never force-pushes.
\texttt{--force-protocol} behaves exactly as described under \texttt{pull}
above.

\subsection{\texttt{autofix}}
\label{sec:autofix}

\begin{lstlisting}[language=bash]
$ pixi run cgitsync autofix [spec.cgs|snapshot.gts] [--error TEXT] [--repo NAME]
\end{lstlisting}

Diagnoses and repairs the situation named by the last failing command's
error, instead of a person reading raw \texttt{git} output and guessing
which of \texttt{merge}/\texttt{pull-force} is safe. With no
\texttt{--error}, reads the most recent \texttt{.cgitsync/logs/*.log}'s
failing command --- the normal workflow is a command fails, run
\texttt{cgitsync autofix} next, and nothing gets retyped. \texttt{--repo}
names which mounted repository to repair when it cannot be guessed from
the error text (a chain-shaped repository's name is usually visible in
its own remote URL, which \texttt{git}'s error text already includes).

\paragraph{Tricky git states, and the safe \cgscli{} command for each.}
The \texttt{SYNC} column \texttt{status} prints (\S\ref{sec:status-sync},
below) names the git-level situation; it cannot say whether the content
behind it is safe to reconcile automatically. This table can:

\begin{center}
\begin{tabular}{lp{8.7cm}}
\textbf{\texttt{SYNC} state} & \textbf{Safe command} \\
\hline
\texttt{ahead(+N)}          & \texttt{push} --- nothing to reconcile, the
  remote is simply behind. \\
\texttt{behind(-N)}         & \texttt{pull} --- nothing to reconcile, this
  workspace is simply behind. \\
\texttt{diverged}, plain-text content (e.g.\ \texttt{.localSpec},
  \texttt{.claude}) & \texttt{merge} by hand, or \texttt{autofix} --- an
  ordinary three-way merge resolves it with nothing lost, since the
  content has no ordering invariant beyond ``these are files''. \\
\texttt{diverged}, chain-shaped content (\texttt{.memory}'s \texttt{lgr/}
  ledger, and --- once built --- \texttt{omniscience}'s register) &
  \texttt{autofix}. \texttt{pull-force} discards whichever side loses the
  race outright; a plain \texttt{merge} leaves an add/add conflict on the
  colliding ledger entry that nothing but \texttt{autofix} knows how to
  resolve without dropping one side's history. \\
Anything \texttt{autofix} does not recognise & Refused outright, with the
  reason printed --- never a guess. \\
\end{tabular}
\end{center}

\texttt{autofix} tells a plain-text divergence apart from a chain-shaped
one before acting on either: fetches, attempts the ordinary merge, and
only on a conflict in a repository registered as chain-shaped does it
re-sequence both sides' new ledger entries --- in the order they were
actually recorded, not ``whichever side pushed first'' --- verify the
result, and complete a real two-parent merge commit. It never
force-pushes and never rewrites history; a repair it cannot verify as
correct is refused, leaving the conflict exactly as \texttt{git} left it,
rather than writing something that looks resolved and is not.

\subsection{\texttt{clone}}

\begin{lstlisting}[language=bash]
$ pixi run cgitsync clone <spec.cgs> [--target-dir DIR] [--output-path DIR]
\end{lstlisting}

Clones the full repository tree declared by the \texttt{.cgs} file. This CLI
command delegates to \texttt{ComplexGitSyncClient.clone}. When
\texttt{--target-dir} is omitted, the project root is derived from the
specification.  \texttt{--output-path} provides the parent directory used to
derive the project root.

For each repository, \cgscli{} tries the repo target branch first and falls
back to the repo fallback branch when the preferred branch does not exist on
the remote. The section ``Branches in a \texttt{.cgs}'' sets out which branch
each repository targets and why. Nested \texttt{.cgs} files are discovered after parent repositories
are cloned, so explicit and automatic nested trees are materialised during the
same command.

\paragraph{What the clone leaves behind.} A clone still downloads one
branch --- that is what keeps it quick. What it no longer does is
\emph{remember} that restriction: it leaves no single-branch fetch refspec
in \texttt{.git/config}, so a later \texttt{fetch} sees the whole remote
rather than the one branch cloned. A workspace cloned before this still
carries the narrow refspec; the first \texttt{pull} or \texttt{push} widens
it, with no re-clone needed. \texttt{status} deliberately does not: it is a
read-only command and never rewrites \texttt{.git/config}.

\subsection{\texttt{checkout}}

\begin{lstlisting}[language=bash]
$ pixi run cgitsync checkout <branch> [--gts <snapshot.gts>] [--ref-kind branch|tag]
\end{lstlisting}

Checks out a branch or tag across the whole tree, parent-first.
Loads the READY registry from \texttt{--gts}, then propagates the ref,
creates missing local branches, and runs \texttt{git checkout} in every repo.
On success the tree remains \texttt{READY} and a new \texttt{.gts} snapshot is written.

\paragraph{A branch that exists on the remote is joined, not recreated.}
When the repository already holds a remote-tracking ref for the branch ---
because some earlier \texttt{pull} downloaded it, or another
\texttt{checkout}/\texttt{branch} fetched it on demand --- \cgscli{} creates
the local branch \emph{from that ref}, with tracking set. You land on your
colleague's commits. When it has no such ref cached, it asks the remote
directly (one \texttt{ls-remote}, no full fetch) before deciding: found, it
fetches that one branch and joins it the same way; genuinely not found, it
creates a new branch where you currently stand, at no extra cost beyond
that one lookup.

This is a change of behaviour, and the reason for it is the failure it
removes. \cgscli{} used to create the branch at the current \texttt{HEAD}
whenever this workspace had not already cached the remote's ref --- so
typing a colleague's branch name, on a clone that had never fetched it,
gave you a branch of that name holding commits that had nothing to do with
theirs, and nothing in the output said so. Depending on some earlier
command (a \texttt{pull}, or a differently-built \texttt{cgitsync} during
\texttt{bootstrap}) to have fetched every branch first was the actual
defect: one command's gap silently became another command's fork.

\texttt{pull} still fetches every branch of each remote before pulling your
own, and remains the cheapest way to bring many refs across at once.
\texttt{checkout} and \texttt{branch} no longer depend on it having run
first: each fetches the one ref it needs, exactly when it needs it.

\subsection{\texttt{close-branch}}

\begin{lstlisting}[language=bash]
$ pixi run cgitsync close-branch <branch> [--private] [--gts <snapshot.gts>]
\end{lstlisting}

Renames \texttt{<branch>} to \texttt{closed/<branch>}, tree-wide,
leaf-first --- locally and on the remote. Requires \texttt{READY}; the
tree remains \texttt{READY} on success. Renames only: nothing is deleted,
so the commits stay exactly as reachable as before, under a name that
says what happened to them. Per repository: the branch is pushed under
its closed name first, the old remote name is removed only once the
commits are reachable under the new one, and only then is the local
branch renamed --- the same order \texttt{memory reboot}'s own branch
archiving already uses. A repository whose remote never had this branch
published is renamed locally only, with nothing to remove remotely.

Refuses, before touching any repository, in two cases: closing the
project's own default branch (every repository with no branch of its own
falls back to it, so removing it would break that fallback), and closing
a branch that any repository in scope is currently checked out on. The
second refusal names which repository and stops there --- this command
closes a branch, it does not move the tree off one; check it out
elsewhere first, then close it. A repository with no local branch named
\texttt{<branch>} is skipped, not an error: a branch that only ever
existed on some repositories (a private/local mount, say) is ordinary,
not a mistake.

\texttt{--private} selects the writable configuration repositories
instead of the project's own, the same as \texttt{branch} (create).

\subsection{\texttt{add}}

\begin{lstlisting}[language=bash]
$ pixi run cgitsync add [PATH ...] [--gts <snapshot.gts>] [--dry-run]
\end{lstlisting}

With no \texttt{PATH} arguments (the default), stages all changes across the
tree (\texttt{git add --all}), leaf-first. Requires \texttt{READY}; tree
stays \texttt{READY}. With \texttt{--dry-run}, \cgscli{} prints the
leaf-first execution plan and does not mutate repositories.

With one or more \texttt{PATH} arguments, each is resolved (relative to
\texttt{CGSHOME}, or absolute) to the one repo in the tree that owns it and
staged there individually (\texttt{git add -- <path>}), leaving every other
repo untouched. A path outside every repo in the tree raises an error
immediately, before anything is staged.

\subsection{\texttt{rm}}

\begin{lstlisting}[language=bash]
$ pixi run cgitsync rm PATH [PATH ...] [--gts <snapshot.gts>] [--dry-run]
\end{lstlisting}

Removes one or more tracked files, each resolved the same way as
\texttt{add}'s targeted form: to the one repo in the tree that owns it,
where it is deleted from disk and the removal staged (\texttt{git rm --
<path>}). A plain tracked file only --- a path that resolves to a directory,
or that does not exist, raises an error rather than partially applying, and
a path outside every repo in the tree raises before anything is removed.
Requires \texttt{READY}; tree stays \texttt{READY}. With \texttt{--dry-run},
\cgscli{} prints the plan without mutating repositories.

Distinct from and unrelated to \texttt{import-submodules}' internal use of
\texttt{git rm --cached} (index-only, dropping a gitlink while keeping the
child's working tree) --- \texttt{rm} deletes the file from disk.

\subsection{\texttt{commit}}

\begin{lstlisting}[language=bash]
$ pixi run cgitsync commit <message> [--gts <snapshot.gts>] [--no-stage] [--dry-run]
\end{lstlisting}

Commits changes across the tree, leaf-first.  Loads the READY registry from
\texttt{--gts}.  Repositories with no staged changes after the optional
\texttt{git add --all} step (suppressed by \texttt{--no-stage}) are silently
skipped. Requires a \texttt{READY} tree; the tree remains \texttt{READY} on success.
Commit preflight now emits warnings for actionable workspace issues such as
dirty worktrees, ahead branches, or stale recorded \texttt{commit\_sha}
values.
With \texttt{--dry-run}, \cgscli{} prints the planned stage/commit sequence
without performing git writes.

The commit message conventionally ends with \texttt{CGS\#VERSION}.

\subsection{\texttt{push}}

\begin{lstlisting}[language=bash]
$ pixi run cgitsync push [--gts <snapshot.gts>] [--dry-run] [--force-protocol {ssh,https}]
\end{lstlisting}

Pushes all repositories to their configured remotes, leaf-first.  Loads the
READY registry from \texttt{--gts}.  Each repo is pushed to
\texttt{entry.remote\_name} (defaulting to \texttt{"origin"}) on its currently
resolved branch.  Refreshes the stored hash in \texttt{GitTree}.  Requires a
\texttt{READY} tree; the tree remains \texttt{READY} on success. Push preflight
blocks detached HEADs, missing remotes, unresolved merges, and
behind/diverged branches; it warns when the local branch is ahead or when
the loaded snapshot records an older \texttt{commit\_sha}.
With \texttt{--dry-run}, \cgscli{} prints the push plan without mutating any
repository. \texttt{--force-protocol} behaves exactly as described under
\texttt{pull} above.

\paragraph{Folds this project's own memory first, when one is mounted.}
Before pushing anything else, \texttt{push} folds whatever has accumulated
in \texttt{.cgitsync} into \texttt{.cgitsync/.memory} and sends it --- the
same work \texttt{cgitsync memory push} does by hand, run automatically the
moment a command was already about to reach a remote. This happens
regardless of \texttt{--private}: a memory, once adopted, settles exactly
as much as an ordinary push. Nothing happens when the tree declares no
memory at all. A memory that is declared but not yet adopted, or one that
cannot reach its own remote for any reason, prints a warning and does not
block the push --- the memory's own trouble is never a reason to fail
somebody else's publish. \texttt{--dry-run} lists \texttt{cgitsync memory
push} first in its plan under the same condition the real run would
perform it.

\subsection{\texttt{freeze-release}}

\begin{lstlisting}[language=bash]
$ pixi run cgitsync freeze-release <name> <message> [--gts <snapshot.gts>] [--dry-run]
    [--force-protocol {ssh,https}]
$ pixi run cgitsync freeze-release-force <name> <message> [--gts <snapshot.gts>] [--dry-run]
    [--force-protocol {ssh,https}]
\end{lstlisting}

Minimalist release workflow from a READY tree. It chains public operations in
order: \texttt{add -> commit -> pull -> push -> freeze}. The force variant uses
\texttt{pull-force} instead of \texttt{pull}. The release name becomes the
shared tag and freeze snapshot label; the message is used for the release
commit.

The \texttt{pull}/\texttt{pull-force} step is skipped (not attempted) when
the current branch has no upstream configured yet — e.g. a branch just
created and checked out this same session, never pushed before. There is
nothing to pull in that case; \texttt{push} (already upstream-aware, adding
\texttt{-u} automatically the same way plain \texttt{git push -u} would)
publishes the branch on its own. On any branch that already has an
upstream, the workflow is unchanged.

\texttt{--force-protocol} is forwarded to the \texttt{pull}/\texttt{pull-force}
and \texttt{push} steps above, behaving exactly as described under
\texttt{pull}; the remote rewrite it makes persists, so the final
\texttt{freeze} step's own tag push picks it up too without needing the
flag itself.

Unlike a bare \texttt{freeze}, this records a \texttt{release} row on the
ledger entry it writes: the installed package's own SemVer and build
counter, plus \texttt{<name>} as the tag actually applied. Name the release
\texttt{v<semver>} (for example \texttt{v3.1.0}) to keep the tag and the
recorded SemVer in agreement.

\subsection{\texttt{freeze}}

\begin{lstlisting}[language=bash]
$ pixi run cgitsync freeze <name> [--gts <snapshot.gts>] [--dry-run]
\end{lstlisting}

Performs a release freeze across all repos, leaf-first: stage/commit,
create+push the shared release tag, refresh SHA/state, and write a
\texttt{.gts} snapshot. Loads the READY registry from \texttt{--gts}.
Requires \texttt{READY} and leaves the tree \texttt{READY}. Freeze preflight
checks remote availability, branch alignment, detached HEADs, unresolved
merges, and tag absence. Dirty worktrees and stale recorded \texttt{commit\_sha} values are
reported as warnings because freeze will stage, commit, and snapshot them.
Freeze snapshots also include a deterministic \texttt{freeze\_manifest} section
that records freeze invariants (immutable snapshot, validated workspace,
synchronized tag, \texttt{release-name}, ledger checkpoint) and declares
\texttt{launch\_state} as the canonical restore operation. Their immutable
snapshot filename includes the release, for example
\texttt{gts-000001-release-2026.05.gts}, while the ledger id remains
\texttt{gts-000001}.
With \texttt{--dry-run}, \cgscli{} prints the planned add/commit/tag/push graph
for the loaded workspace without mutating it.

\paragraph{Folds this project's own memory first, same as \texttt{push}.}
\texttt{freeze} (and, through it, \texttt{tag}, which folds the same way on
its own combined commit-tag-push) crosses the \texttt{.cgitsync}
$\rightarrow$ \texttt{.cgitsync/.memory} frontier before doing anything
else, under the same conditions \texttt{push} does. \texttt{freeze-release}
folds twice in one run --- once via its own \texttt{push} step, once via
its own \texttt{freeze} step --- which is harmless: a fold with nothing new
pending commits nothing and only re-sends the branch.

\subsection{\texttt{import-submodules}}

\begin{lstlisting}[language=bash]
$ pixi run cgitsync import-submodules REPO_ROOT [--apply] [--recursive]
\end{lstlisting}

Reads \texttt{REPO\_ROOT/.gitmodules} and reports (or converts) every declared
git submodule to a plain ComplexGitSync nested clone. This is the whole job:
turning gitlinks into plain clones on disk. It does not author a
\texttt{.cgs} --- \texttt{.gitmodules} never records the root's own identity,
and a checkout worth importing already has a \texttt{.cgs} (hand-authored)
or can get one from \texttt{discover} (\S\ref{chap:user-guide}
``\texttt{discover}''), before or after applying.

Without \texttt{--apply} (the default), the command performs a
\emph{dry run}: it prints each submodule's name, path, URL, and branch,
without touching the repository.

\texttt{--recursive} (same name and meaning as \texttt{git submodule
update}'s own flag) also converts any submodule that itself has its own
checked-out \texttt{.gitmodules}, at any depth --- root first, not leaf
first: converting a submodule stages changes in its own working tree, and
converting a deeper level first would make the shallower level's own
preflight check (is the submodule's working tree clean?) see that staged
dirt and reject its own conversion; converting the shallower level first
avoids this, since \texttt{git rm --cached <path>} only ever touches the
parent's own index, never the submodule's working tree. A submodule path
that was never checked out (\texttt{git submodule update --init} not run
for it) has no \texttt{.git} to read a nested \texttt{.gitmodules} from and
is invisible either way --- the same ceiling \texttt{discover} already has
for itself. Without \texttt{--recursive}, behavior is unchanged: exactly
\texttt{REPO\_ROOT/.gitmodules}, one level.

The report lists every level together. Each entry's \texttt{path} is
counted from \texttt{REPO\_ROOT}, and \texttt{declared in} names the
\texttt{.gitmodules} file it was read from. A submodule's own
\texttt{.gitmodules} path means nothing on its own once more than one
repository is involved --- in \texttt{cawaqsviz}, two of the three
submodules are called \texttt{docs/...} by their own repository.

With \texttt{--apply}:
\begin{itemize}
  \item Verifies every submodule's working tree is clean (\texttt{git status
        --porcelain} must be empty apart from that repository's own
        \texttt{.gitignore}); rejects dirty trees immediately.  The
        \texttt{.gitignore} exemption exists because \cgscli{} writes that
        one file itself in every repository holding a child, so it is
        routinely dirty in exactly the tree being converted; no conversion
        can harm it, since \texttt{git rm --cached} never touches the
        child's working tree.
  \item Runs \texttt{git rm --cached <path>} per submodule, dropping the
        gitlink from the index while leaving the child's working tree and
        \texttt{.git} directory intact (no re-clone, no local history lost).
  \item Removes the converted stanza from \texttt{.gitmodules} (deletes the
        file entirely when all stanzas are converted) and stages the change.
  \item Appends \texttt{<path>} to the parent's \texttt{.gitignore} using the
        same helper that the \texttt{.gitignore} lifecycle sync uses.
\end{itemize}

After applying, review and commit the staged changes manually:
\begin{lstlisting}[language=bash]
$ git -C REPO_ROOT status          # deleted .gitmodules, modified .gitignore
$ git -C REPO_ROOT commit -m "chore: retire git submodules"
\end{lstlisting}

\textbf{Note:} ComplexGitSync stores no credentials and has no private-repo
authentication story.  If a submodule's upstream is private and the ambient
environment lacks access, the subsequent \cgscli{} \texttt{initialise} or
\texttt{pull} will fail at the clone step.  The \texttt{import-submodules}
step itself (which only touches the local index and working tree) will still
succeed.

\subsection{\texttt{init-from-submodules}}

\begin{lstlisting}[language=bash]
$ pixi run cgitsync init-from-submodules REPO_ROOT [--cgs FILE]
      [--max-depth N] [--dry-run] [--force] [--force-protocol {ssh,https}]
\end{lstlisting}

Adopts a submodule-based checkout in one command: \texttt{discover} the
repositories under \texttt{REPO\_ROOT}, write the drafted \texttt{.cgs},
\texttt{initialise} the tree in place, then convert every submodule with
\texttt{import-submodules --recursive --apply}.  \texttt{REPO\_ROOT} is a
project cloned and \texttt{git submodule update --init --recursive}'d by
hand; the command performs no clone of the root itself.

\textbf{The conversion must come last, and this command is the reason
why.}  \texttt{initialise} adopts the root in place but deletes and
re-clones every \emph{other} repository straight from its remote, and those
remotes still declare submodules --- so a conversion run first is undone
for every repository except the root.  Nor can a second pass repair that:
\texttt{import-submodules --recursive} follows the submodule graph declared
by the root's own \texttt{.gitmodules}, which the first pass removed, so no
deeper level stays reachable.  Running the two commands in the other order
by hand produces a half-converted tree that looks finished.

Three conditions are checked before anything is written, so a run that
cannot work fails while the checkout is still untouched:

\begin{itemize}
  \item the discovery resolved at least one repository (skipped when
        \texttt{--cgs} supplies a file, which is then the authority);
  \item the project name matches the directory name --- \texttt{CGSHOME} is
        resolved as \texttt{<parent>/<project-name>}, so a mismatch would
        send \texttt{initialise} to a sibling directory that does not
        exist;
  \item \texttt{REPO\_ROOT} has a \texttt{.gitmodules}.  Without one there
        is nothing to convert, while the clone step would still delete and
        re-clone every non-root repository, discarding any uncommitted work
        in them.  \texttt{--force} overrides this one.
\end{itemize}

\texttt{--dry-run} prints the discovery and the submodules that would be
converted without writing a \texttt{.cgs}, cloning, or touching any
repository.  \texttt{--cgs FILE} uses an existing \texttt{.cgs} instead of
writing one (or writes the draft there, when \texttt{FILE} does not exist).
\texttt{--max-depth N} bounds the discovery step, which is otherwise
unbounded --- most trees never need it.  \texttt{--force-protocol} applies
to the clone step only; discovery reads each repository's configured
remote as it is.

The command stops at a \texttt{READY} tree with the conversion staged but
not committed --- it touches every repository that held a submodule, and
some of those may belong to other people --- and prints the
\texttt{branch}/\texttt{checkout}/\texttt{add}/\texttt{commit} sequence to
run next.

\subsection{\texttt{memory}}

\begin{lstlisting}[language=bash]
$ pixi run cgitsync memory status [--search-dir DIR]
$ pixi run cgitsync memory list   [--search-dir DIR]
$ pixi run cgitsync memory show <state> [--full] [--search-dir DIR]
$ pixi run cgitsync memory explore [--branch NAME] [--timeline] [--search-dir DIR]
\end{lstlisting}

Looks at what this workspace remembers. A \textbf{State} is one recorded
snapshot of the tree, named by its own content; the \textbf{ledger} holds
one entry per operation, saying when that State was seen and by which tool
versions.

\begin{itemize}
  \item \texttt{status} --- how many States, how long the chain is, when it
        was last written, which of \texttt{verify}'s four answers this
        memory is in, and the toolchain the records carry. The first
        entry's versions are shown beside the latest when they differ, so a
        chain spanning an upgrade says where the upgrade fell.
  \item \texttt{list} --- every State, newest recording first, with the
        commands that recorded it. A State with no timestamp is one no
        entry records: history from before the ledger, or a file that
        arrived some other way.
  \item \texttt{show <state>} --- the environment that State's own commits
        ran under, as a bare reference; that State's own topology right
        after, drawn exactly the way \texttt{view-tree} draws the live one
        (\texttt{format\_view\_tree}) but as it was at that State, not as it
        is now; then every entry naming it and what each \texttt{commit}
        wrote: the message, the repository, and whether the commit has been
        published. The argument is the State's content hash or any
        unambiguous prefix of it; nobody retypes sixty-four characters.
        Pasted back verbatim from wherever it was printed also works: the
        bare prefix, the \texttt{state=<hash>} label \texttt{explore
        --timeline} prints it under, or the full \texttt{state(<hash>)} form
        a ledger entry uses. \texttt{--full} prints whole messages instead
        of first lines, and the remote and ref each publication went to.
  \item \texttt{show env=<ref>} --- one Environment record, in full: machine,
        tools, credentials and manifest digests, drawn as the same kind of
        tree \texttt{view-tree} draws a repo tree with. This is the detail a
        State's own bare \texttt{environment=env(<hash>)} reference never
        shows, since a State answers "what was this tree", not "what ran
        it". Takes the full id, a bare prefix, or the \texttt{env=<hash>}
        form its own argument is spelled with.
  \item \texttt{explore} --- a memory a person can read without a hash.
        The default view is one row per commit this memory recorded as
        published, newest push first, on the branch checked out here ---
        what a colleague pulling this branch would see. \texttt{--timeline}
        reads the ledger straight through instead: one row per entry, in
        the order it happened, so \texttt{checkout}/\texttt{merge}/
        \texttt{push} appear beside \texttt{commit} rather than being
        dropped --- each row also prints the State it recorded, so a prefix
        to hand \texttt{show} is never more than a scroll away.
        \texttt{--branch NAME} asks for a different memory branch;
        only the one actually checked out here can be read this way today,
        and a mismatch names \texttt{memory clone --branch NAME} rather
        than answering for the wrong branch.
\end{itemize}

A commit marked \texttt{unpushed} exists on this machine and nowhere else.
That distinction is the reason the messages are kept at all: while the
repository is present \texttt{git log} is the better answer, and this record
is for when it is not --- a deleted branch, an archived project, a fork that
went away. The ledger entry that recorded a set of messages carries their
fingerprint, so editing one afterwards is a finding \texttt{verify} reports.

\texttt{status}, \texttt{list}, \texttt{show} and \texttt{explore} are
read-only; nothing they do sends a memory anywhere.

\paragraph{Keeping a memory when the disk does
not.}\label{sec:memory-mount} A memory can be a
repository of its own, mounted at \texttt{.cgitsync/.memory} exactly as
\texttt{.localSpec} is mounted: one repository holding every project's
memory, on a branch per project. \texttt{.cgitsync} itself stays the
workspace's own live state --- States, the ledger, logs --- so the mount
one level inside it is checked out and merged like any other private
repository, never blocked by activity that is not its own. Getting there
is five commands, run once per project; \texttt{tutorials/05\_memory.md}
walks all five in order.

\begin{lstlisting}[language=bash]
$ pixi run cgitsync repo create github:you/.memory
$ pixi run cgitsync memory mount [--cgs FILE]
$ pixi run cgitsync memory adopt
$ pixi run cgitsync memory push
$ pixi run cgitsync memory branch --project-branch main
\end{lstlisting}

\begin{itemize}
  \item \texttt{repo create} --- creates the repository by running the
        provider's own tool (\texttt{gh}, \texttt{glab}, \texttt{tea}),
        which you have already signed in to. ComplexGitSync stores no
        credential and sends none. A repository that already exists is an
        ordinary success; a tool that is missing or signed out prints the
        command to run and exits \texttt{2}.
  \item \texttt{memory mount} --- adds one entry to a \texttt{.cgs} that
        already exists, leaving the rest of the file, comments included,
        exactly as it was. Adding it twice changes nothing.
  \item \texttt{memory adopt} --- creates \texttt{.cgitsync/.memory} fresh
        and turns it into that repository, on this project branch's memory
        branch. It starts empty: the first \texttt{memory push} is what
        folds in whatever \texttt{.cgitsync} already held. \texttt{--reboot}
        adopts the same repository identity but skips carrying forward
        whatever the fallback branch already holds, so the mount's content
        starts fresh instead (\S\ref{sec:memory-reboot}).
  \item \texttt{memory branch} --- creates the memory branch another
        project branch will need, and pushes it. Needed once, for a memory
        born on a feature branch: merging \texttt{memory-dev} into
        \texttt{main} asks for a branch that has never existed. A merge
        that finds one missing names this command rather than creating the
        branch itself, because a branch that is not there is as likely to
        be a typing mistake as a new branch.
\end{itemize}

\paragraph{Migrating a memory mounted before \texttt{memory migrate}
existed.} A memory adopted before this command existed sits directly at
\texttt{.cgitsync}, sharing it with the workspace's own live state.
\texttt{cgitsync memory migrate [--cgs FILE]} moves it onto the layout
above --- \texttt{.git} and every tracked file travel down into
\texttt{.cgitsync/.memory}, history unchanged; whatever was never tracked
stays exactly where it already was, which is where the new layout wants
it anyway --- and updates the \texttt{.cgs} entry that declares the mount
to match. A fresh \texttt{memory adopt} never needs this: it mounts at
the new layout directly.

\begin{lstlisting}[language=bash]
$ pixi run cgitsync memory init    # the .cgs entry to add, and the branch
$ pixi run cgitsync memory clone   # bring a memory onto a new machine
$ pixi run cgitsync memory push    # commit what it gained, and push it
\end{lstlisting}

\texttt{init} proposes and stops: ComplexGitSync speaks Git and nothing
else, so it prints the command that creates the repository and waits.
Nothing is ever pushed automatically --- a machine with no network keeps a
complete, verifiable memory and sends it later.

\texttt{clone} refuses to write over a memory that is already here, because
a memory nobody has pushed is the only copy of itself.

\paragraph{Starting a memory's history over.}\label{sec:memory-reboot} A
project's shape changes --- repositories added, removed, restructured ---
and a memory built for the old shape stops being a clean answer to "what
does this project look like." \texttt{cgitsync memory reboot} closes the
current chapter and opens an empty one, without losing the old one:

\begin{lstlisting}[language=bash]
$ pixi run cgitsync memory reboot
\end{lstlisting}

In order: whatever \texttt{.cgitsync} is holding pending is folded in and
pushed under the branch's current name, exactly as \texttt{memory push}
already does; the tree's current shape is exported --- from the loaded
\texttt{.gts}, never a hand-authored file --- to a permanent, versioned
\texttt{.cgitsync/.memory/.cgs/<project>-v<N>.cgs} (\texttt{N} incrementing
once per reboot, never overwritten, never reused); the branch is archived,
renamed both locally and on origin to \texttt{<branch>.archived-<date>},
the old name removed from origin only \emph{after} the new one already
holds the same commits, so nothing the memory ever recorded is
unreachable for even a moment; and a fresh, empty branch takes the
original name, so nothing about how the memory is mounted changes.
Nothing is ever force-pushed or deleted. \texttt{cgitsync memory adopt
--reboot} is the same fresh start for a mount being adopted for the very
first time. Appending --- the ordinary \texttt{memory adopt}, and the only
thing \texttt{memory push} ever does --- stays the default either way.

\paragraph{What a memory carries off your machine.} One path: the tree's own
root, with \texttt{\$HOME} substituted. Every other path it records --- each
repository's, and any path in a command line it logged --- is written
against the tree as \texttt{\$CGSTREE/...}, so neither your directory
layout nor your user name travels with it.

\subsection{\texttt{self-history}}

\begin{lstlisting}[language=bash]
$ pixi run cgitsync self-history add \
    --ticket NAME --goal TEXT --action TEXT \
    --worker-role ROLE --worker-vendor NAME --worker-model NAME \
    --orchestrator-role ROLE --orchestrator-vendor NAME --orchestrator-model NAME \
    --spec-respect-score N --spec-respect-basis measured|asserted --spec-respect-reasoning TEXT \
    --gating-score N --gating-basis measured|asserted --gating-reasoning TEXT \
    --quality-score N --quality-basis measured|asserted --quality-reasoning TEXT \
    [--state-before STATE] [--state-after STATE] \
    [--lint-passed|--lint-failed] [--tests-passed|--tests-failed] \
    [--pushed] [--pushed-reason TEXT] [--search-dir DIR]
\end{lstlisting}

Records one piece of agent work: which ticket it served, the goal and the
action taken in plain English (at most three lines each), the two agents
involved --- a \textbf{worker} that implemented the ticket and an
independent \textbf{orchestrator} that quoted the work, each with a role
from \texttt{.localSpec/AGENT.md}'s roster, a vendor and a model version
--- and a three-part conformity score (spec respect, \texttt{.PUBLIC}/
\texttt{.PRIVATE} gating, quality of production), each part marked
\texttt{measured} or \texttt{asserted} rather than blending the two.

Two facts the command fills in itself, because this project can check them
directly rather than trust an agent to type them: which signed
agent-contract record was in force (by its own hash, from
\texttt{.agent/.distant/dev-sync/agent-contracts/current}), and this
workspace's own \texttt{errors=} count from the same view \texttt{status}
prints. Everything else is the orchestrator's own account. A missing
signed contract is not an error --- the field is simply empty.

\texttt{--state-before}/\texttt{--state-after} are \textbf{verified against
the ledger}, not merely shape-checked: each, if given, must name a State
some ledger entry actually recorded, still on disk, still hashing to its
own name --- the same three questions \texttt{verify} itself asks. A
citation that fails any of them is refused, not recorded as though it were
fact. \texttt{--state-after} defaults to the ledger's own most recent
entry when not given.

The record is written to \texttt{.cgitsync/.self-history/}, content-addressed
the same way a State or an Environment record is --- the pending half.
Once self-history has been adopted (below), \texttt{cgitsync memory push}
folds it into \texttt{.cgitsync/.memory/.self-history/} and pushes it,
\emph{before} folding and pushing \texttt{.memory} itself, so the parent's
own push always records a leaf that already exists.

\begin{lstlisting}[language=bash]
$ pixi run cgitsync self-history adopt [--owner NAME] [--branch NAME] [--search-dir DIR]
\end{lstlisting}

Self-history rides \texttt{.memory}'s own lifecycle rather than needing a
lifecycle of its own. \texttt{self-history adopt} is the \emph{one} place
this is ever decided: it writes \texttt{config-memory.cgs} once, for a
brand-new project or as the explicit retrofit for a \texttt{.memory}
adopted \emph{before} self-history existed. From then on,
\texttt{cgitsync memory adopt} follows that decision automatically,
silently, whenever a project's own memory branch inherits it from a
fallback base branch that already carries it --- the same mechanism that
already lets a new project's memory share history with another. A project
that has never adopted self-history has no such file anywhere, so this
sees no new behaviour at all. Unlike that automatic follow, asking for
\texttt{self-history adopt} explicitly raises when the remote is not
there, rather than quietly doing nothing --- asking for it means wanting
to know why it did not work.

\texttt{cgitsync memory clone} brings self-history back too, the moment it
was adopted upstream; \texttt{cgitsync memory reboot} leaves it entirely
alone, on purpose --- a reboot starts \texttt{.memory}'s own history over,
and the record of what happened, including across the reboot, is exactly
the thing that should survive it.

\begin{lstlisting}[language=bash]
$ pixi run cgitsync memory self-history [--search-dir DIR]
\end{lstlisting}

Every self-history record this workspace holds, folded and pending,
oldest first --- consultation only. A workspace that has never adopted
self-history, or has adopted it but written nothing yet, prints that
plainly rather than erroring.

\subsection{\texttt{verify}}

\begin{lstlisting}[language=bash]
$ pixi run cgitsync verify [--search-dir DIR] [--repair]
\end{lstlisting}

Says which of five things is true of this workspace's recorded history.
One of them, never a blur of two:

\begin{center}
\begin{tabular}{llp{7.6cm}}
\textbf{Answer} & \textbf{Exit} & \textbf{Means} \\
\hline
\texttt{verified}   & \texttt{0} & A chain was read and every link held. \\
\texttt{no-history} & \texttt{0} & Nothing has been recorded here yet. Not
                                    a failure --- a new workspace is not a
                                    broken one. \\
\texttt{legacy}     & \texttt{1} & History exists, in the single-file
                                    \texttt{.lgr} register, which carries no
                                    chain. Readable, not verifiable: an edit
                                    to it leaves no trace. \\
\texttt{corrupt}    & \texttt{1} & A chain was read and it does not hold. \\
\texttt{time-inconsistent} & \texttt{1} & The chain held, but its own
                                    timestamps move backwards somewhere.
                                    Your history is intact; the clock that
                                    stamped it was not. \\
\end{tabular}
\end{center}

\texttt{legacy} exits non-zero on purpose. The question is ``is this
history intact?'', and ``I cannot tell'' is not a yes: a build that gates on
\texttt{verify} must not pass because the evidence is in a format nobody can
check. Every workspace created before the hash-chained register is written
is in exactly that state.

\texttt{time-inconsistent} is deliberately not \texttt{corrupt}. Every link
checked out --- nothing was rewritten --- so the two answers send you to
look at different things. A clock corrected mid-session, a restored virtual
machine, or a machine that disagreed about the hour all read like this, and
so does an entry somebody backdated: a date cannot be moved backwards
without contradicting the chain around it. It still exits non-zero, because
something is wrong even though your history is not.

When there \emph{is} a chain, what is checked: every entry's \texttt{prev}
must match the previous entry's \texttt{entry\_hash}
(\texttt{BROKEN\_LINK}), every entry's own \texttt{entry\_hash} must match
its recomputed canonical hash (\texttt{BAD\_ENTRY\_HASH}), \texttt{seq}
numbers must be gap-free and unique
(\texttt{SEQ\_GAP}/\texttt{SEQ\_DUPLICATE}), and the cached \texttt{HEAD}
must agree with the recomputed head (\texttt{HEAD\_STALE}).

\texttt{CGSHOME} is located the same way as \texttt{status}: from
\texttt{--search-dir}, else \texttt{\$CGSHOME}, else the current working
directory's ancestors.

With \texttt{--repair}, a stale \texttt{HEAD} cache is corrected to match the
recomputed true head. \textbf{Entries themselves are never rewritten or
deleted} by \texttt{verify}, repair or not — a broken chain is reported, not
silently healed, since a register that can be edited back to looking clean
provides no evidence of anything.

\paragraph{The States on disk are checked too.} Beyond the chain itself,
\texttt{verify} cross-references what the entries say against what is
there: an entry naming a State that is missing (\texttt{MISSING\_STATE}), a
stored snapshot whose content no longer hashes to the name it is filed
under (\texttt{STATE\_DIGEST\_MISMATCH}), and a State on disk that no entry
records (\texttt{ORPHAN\_STATE}). The second is what naming a State by its
content bought: a stored snapshot can no longer be edited quietly.

\paragraph{Where the history lives.} \texttt{.cgitsync/lgr/}, one file per
entry, appended to by every command that writes a State. Each entry records
the versions that produced it --- cgitsync, git, pixi, and the data backends
dvc and git-lfs --- inside the entry hash, so an edited version string is as
detectable as an edited command. The older single-file \texttt{.lgr}
register is still read, and is no longer written.

\subsection{\texttt{env} and \texttt{env check}}

\begin{lstlisting}[language=bash]
$ pixi run cgitsync env [--search-dir DIR]
$ pixi run cgitsync env check [--search-dir DIR] [--cgs FILE]
\end{lstlisting}

\texttt{env} observes the environment that can make the current tree usable.
It reports the machine architecture, Pixi platform, operating system, kernel
and libc; raw and parsed versions for Python and the tools \cgspkg{} drives;
provider authentication as a fact; and the paths and digests of recognised
manifests in the tree. A missing tool is printed as \texttt{none}, never as
an empty value. Tokens, keys, user names, key paths, and paths outside
\texttt{\$CGSTREE} are never part of an Environment record.

\texttt{env check} compares that observation with requirements from the
loaded tree's source \texttt{.cgs}. Pass \texttt{--cgs FILE} to check an
explicit authoring document instead. Missing tools and versions older than
the declared minimum are drift: the command names each mismatch and exits
non-zero, which makes it suitable for continuous integration. Extra observed
tools are reported as information and do not make the check fail.

The check is separate from tree operations on purpose. \texttt{initialise}
and \texttt{pull} may warn about drift, but a mismatch never blocks them.
Both environment commands are read-only and install nothing.

The optional authoring declaration names the repository where the environment
is installed and requirements that cannot be inferred from a lock file:

\begin{lstlisting}
environment_root = "my-project"

[environment]
tools = { git = "2.43", pixi = "0.66" }
compilers = ["cc"]
system_libraries = ["libssl"]
services = ["postgresql"]
manifests = ["Cargo.lock"]
\end{lstlisting}

The root must name a configured repository (or \texttt{root}). Manifest
contents remain in Git; the Environment record stores only tree-relative
paths and digests.

\subsection{\texttt{launch-release}}

\begin{lstlisting}[language=bash]
$ pixi run cgitsync launch-release <name> [--gts <snapshot.gts>]
\end{lstlisting}

Checks out the frozen release tag \texttt{<name>} across the loaded READY tree,
parent-first, and writes a fresh \texttt{.gts} snapshot.  This is the release
oriented CLI path for returning a workspace to a frozen version; direct tag
creation remains available from the Python client API as \texttt{tag()}.

The integration suite validates local no-network lifecycle coverage through
the CLI path, including initialisation and the READY git command cycle
\texttt{add -> commit -> push -> freeze -> launch-release}.

\subsection{\texttt{view-tree}}

\begin{lstlisting}[language=bash]
$ pixi run cgitsync view-tree [spec.cgs|snapshot.gts] [--depth N] [--collapse REPO]
\end{lstlisting}

Renders a deterministic terminal tree that focuses on repository topology and
operational state. Each line reads:

\begin{center}
\texttt{name (node\_type) [sync\_state] @sha br=<branch> [fb=<fallback>]}
\end{center}

\texttt{br=} is printed for \emph{every} repository, including the common
case where the branch is simply the tree's default, so that reading
\texttt{view-tree} tells you the branch each repository targets without
opening the \texttt{.cgs}. \texttt{fb=} appears only when the declared
fallback branch differs from that target, since a fallback equal to the
target says nothing. Use \texttt{--depth} to limit visible levels and repeat
\texttt{--collapse} to hide selected subtrees.

On a tree with private mounts, \texttt{br=} is where you see pinning at work:
the unprivate repositories carry the branch you moved the tree to, and each
private mount still shows its own.

\textbf{Branch Topology Propagation Rules (T35):}

\begin{enumerate}
  \item \textbf{Reference branch}: The root repository's current branch is
        the canonical reference for all repos.
  \item \textbf{Leaf-to-root inheritance}: Branch targeting flows
        root-first via \texttt{propagate\_global\_branch}, verified
        on-disk by \texttt{ComplexGitSyncClient.validate\_topology()}
        (Python API only) without issuing any git writes.
  \item \textbf{Allowed divergence}: Repositories in a frozen
        (\texttt{TAG}) state produce a \texttt{tag\_divergence} diagnostic
        that does \emph{not} make the topology incoherent.
  \item \textbf{Incoherent states} (blocking): \texttt{misaligned\_branch}
        (different branch from root) and \texttt{detached\_head} (detached
        HEAD without a tag reference).
\end{enumerate}

\subsection{What \texttt{cgitsync} promises a script}

\paragraph{Exit codes.} Three, meaning the same thing in every command:

\begin{center}
\begin{tabular}{lp{10.5cm}}
\textbf{Code} & \textbf{Means} \\
\hline
\texttt{0} & The command did what was asked. \\
\texttt{1} & It ran, and the answer is no --- a merge conflict, a tree that
             is not \texttt{READY}, a verification that found something. \\
\texttt{2} & It could not run --- bad arguments, no workspace, a missing or
             unreadable file. \\
\end{tabular}
\end{center}

The distinction that matters to a script is between ``I asked and the answer
is no'' and ``I could not ask''. \texttt{validate} is the one command that
judges a document rather than acting on one, so an invalid document is its
answer (\texttt{1}); every other command exits \texttt{2} on the same
document because it could not run at all.

An expected failure prints one line on stderr and no traceback. A traceback
is therefore a defect report, not a usage error --- which is the point of
mapping the expected ones.

\paragraph{\texttt{--json}.} \texttt{status} and \texttt{verify} accept
it. Each prints one JSON object on stdout and nothing else; every
human-facing line goes to stderr instead, and the exit code is unchanged, so
a caller may use either signal. A failure prints one object too, carrying
\texttt{"status": "error"} and the exit code, rather than leaving a pipe
with nothing to parse. \texttt{--help}, \texttt{--version} and a command
line that fails to parse are deliberately not JSON: the flag itself may be
what failed.

The object carries \texttt{schema\_version}. The promise attached to it is
\textbf{additive only}: new fields may appear, existing fields do not change
meaning and do not vanish.

\subsection{Where a command works: the default workspace}

Every command that is not given an explicit path finds its workspace from
three inputs, in this order: \texttt{--search-dir}, then
\texttt{\$CGSHOME}, then the ancestors of the current directory. When none
of them lands on a \texttt{.cgitsync} directory there is a fourth answer: an
empty workspace of ComplexGitSync's own, under \texttt{\$HOME/.cgs} (or
\texttt{\$CGSPATH} when that is set).

That workspace is a real directory holding a valid \texttt{.gts} with no
repositories in it. An empty tree is not an error --- it is a project that
has not started --- so \texttt{status} answers \texttt{no living project
yet}, names the three commands that start a project, and exits \texttt{0}.

\begin{itemize}
  \item \textbf{Created once, then reused.} Its path is recorded in
        \texttt{\$HOME/.cgs/default}; later runs read that pointer instead
        of minting a second empty workspace.
  \item \textbf{It never chooses for you.} When other workspaces exist
        under the root they are listed, each with the
        \texttt{export CGSHOME=...} line that would select it, and none is
        selected automatically.
  \item \textbf{It never overrides \texttt{--search-dir}.} A directory you
        named that holds no workspace is an error, not an invitation to work
        somewhere else.
\end{itemize}

\paragraph{The \texttt{use\_case} field.} Every discovery line also says
which of the two ways of running is in force: \texttt{nested} when the
ComplexGitSync being executed lives inside the workspace it manages,
\texttt{standalone} otherwise. It is observed and printed, never obeyed ---
nothing behaves differently because of it. It exists so that a user who
believes they are in one case while standing in the other is told.

\subsection{\texttt{status}}

\begin{lstlisting}[language=bash]
$ pixi run cgitsync status [--gts FILE]
\end{lstlisting}

Summarises tree readiness and per-repository sync state. The table reports
the scope, local branch, upstream branch, local worktree state,
ahead/behind tracking counts, current HEAD, and the commit recorded in the
loaded \texttt{.gts}.

\paragraph{The \texttt{cgitsync\_branch} field.} The \texttt{summary} line
carries \texttt{cgitsync\_branch=<branch>}: the branch the project is on,
which is the branch its root repository is on. It is what
\texttt{cgitsync branch} and \texttt{cgitsync checkout} set, and what every
other repository follows.

The table shows a branch per repository and they are not all the same on
purpose --- a \texttt{private/local} repository keeps this project's
settings on a branch named after the project, so a project on
\texttt{apoub} shows \texttt{ComplexGitSync\_apoub} there. The field is the
one place that says which of those branches is the project's own.

\begin{center}
\begin{tabular}{lp{10.5cm}}
\textbf{Shown} & \textbf{Means} \\
\hline
\texttt{<branch>}  & The branch the root repository is on. \\
\texttt{detached}  & The root is parked on a commit rather than a branch.
                     \texttt{cgitsync checkout <branch>} puts the tree
                     back. \\
\texttt{unknown}   & There is no branch to report: no project has been
                     loaded, or Git could not be asked. \\
\end{tabular}
\end{center}

The field is printed in every state. A field that disappears when the
answer is awkward is harder to read, by eye or by script, than one that
says it does not know.

\paragraph{The \texttt{SCOPE} column.} It says, in three words, what kind
of repository each row is and whether you may write to it:

\begin{center}
\begin{tabular}{lll}
\textbf{Shown} & \textbf{Means} & \textbf{In the \texttt{.cgs}} \\
\hline
\texttt{project}         & this project's own; write freely       & no \texttt{private} \\
\texttt{private/local}   & shared, and yours to write             & \texttt{private, writable} \\
\texttt{private/distant} & shared, read-only                      & \texttt{private} alone \\
\end{tabular}
\end{center}

\textbf{private} means the repository is shared with other projects --- a
configuration repository. \textbf{local} then means this project may write
to it, normally because it sits on a branch named after your project;
\textbf{distant} means you only read it. A repository nested inside a
private one is shown the same way, because it is just as shared.

\texttt{commit} and \texttt{push} reach the \texttt{project} rows;
\texttt{--private} reaches the \texttt{private/local} rows; nothing writes
to a \texttt{private/distant} row. When any private row is present, the
table prints a legend saying so.

\paragraph{The \texttt{SYNC} column.}
\label{sec:status-sync}
It says how the branch a repository
is on stands against the branch it tracks --- the one named in
\texttt{UPSTREAM\_BRANCH}, which reads \texttt{-} when there is none. There
are six values:

\begin{center}
\begin{tabular}{lp{10.5cm}}
\textbf{Shown} & \textbf{Means} \\
\hline
\texttt{synced}            & Level with the upstream. \\
\texttt{ahead(+N)}         & \texttt{N} commits here that the remote does
                             not have. \texttt{push} sends them. \\
\texttt{behind(-N)}        & \texttt{N} commits on the remote that are not
                             here. \texttt{pull} fetches them. \\
\texttt{diverged(+N/-M)}   & Both, from a common ancestor. \texttt{autofix}
                             (\S\ref{sec:autofix}) diagnoses which of
                             \texttt{merge} or a chain-aware splice is
                             safe; \texttt{pull-force} discards this
                             workspace's own \texttt{N} commits
                             unconditionally and should not be the first
                             thing tried. \\
\texttt{no-upstream}       & This branch was never pushed, so there is
                             nothing to compare it to. Normal for a branch
                             you just made, and for a \texttt{private/local}
                             repository that only \texttt{push --private}
                             ever sends. \\
\texttt{unknown}           & The branch names an upstream that does not
                             resolve. \texttt{pull} or \texttt{push} repairs
                             it; if it persists, the remote is unreachable
                             or the ref was deleted. \\
\end{tabular}
\end{center}

The last two are worth telling apart. \texttt{no-upstream} is a normal
state of a branch nobody has pushed yet; \texttt{unknown} is a fault. They
used to share the one word \texttt{unknown}, which made the column not
worth reading --- an ordinary new branch looked exactly like a broken one.

\paragraph{\texttt{unmeasured} in the \texttt{summary} line.} The summary
counts the \texttt{no-upstream} and \texttt{unknown} rows as
\texttt{unmeasured}, separately from \texttt{ahead} and \texttt{behind}.
Those rows were previously folded into \texttt{ahead=0 behind=0}, which
reported a measurement nobody had taken. A repository nobody could measure
is not a repository that is level, and the summary no longer says the
second when it means the first.

\paragraph{The memory mount reading \texttt{dirty} mid-\texttt{push}, or
\texttt{HEAD} ending in \texttt{*} right after, is normal.} If the tree
mounts its own memory at \texttt{.cgitsync/.memory}
(\S\ref{sec:memory-mount}), that row is an ordinary private/local
repository the rest of the time --- clean before and after every command
except \texttt{memory push} itself, which folds \texttt{.cgitsync}'s own
pending content into it before committing. One thing about it is still
worth knowing:

\begin{itemize}
  \item \textbf{\texttt{HEAD} ends in \texttt{*} right after
        \texttt{memory push}.} The \texttt{*} means the loaded \texttt{.gts}
        recorded a different commit than the one actually there now.
        \texttt{memory push} moves the memory's real commit directly and
        does not write a new \texttt{.gts} itself, so the State most
        recently written still names the commit from before the push. The
        very next command that writes a State reads the memory's true
        \texttt{HEAD} and the two agree again --- it never needs a second
        \texttt{memory push} to catch up, only an ordinary next command.
\end{itemize}

Not worth chasing to a clean reading on purpose: a memory that never
showed this for a moment would not be recording what \texttt{memory push}
just did.

\subsection{\texttt{merge}}

\begin{lstlisting}[language=bash]
$ pixi run cgitsync merge <branch> [--into TARGET] [--private] [--ff-only|--no-ff] [--dry-run] [--resolve]
\end{lstlisting}

Merges a branch across the tree, leaf-first. The argument is always the
\textbf{project's} branch; each repository works out what that means for
itself. A project-owned repository merges that branch. A private/local
repository merges the branch derived from it
(\texttt{<base>\_<branch>}), because that is where its settings for that
project branch live --- it has no branch of the project's own name at all.

\paragraph{\texttt{--into}: merging into a branch you are not on.} Without
it, \texttt{merge} merges into whatever is checked out, so reaching another
branch means running \texttt{checkout} first. \texttt{--into TARGET} names
the target and does both halves in one command.

That is not a convenience when the tree being merged \textbf{contains the
ComplexGitSync you are running} --- the developer checkout, which installs
itself editable. A separate \texttt{checkout} replaces that build, and the
merge typed after it runs under the older one, against a workspace the newer
one wrote. A single command finishes under the build it started with,
because Python has already imported its modules.

Every repository is checked before any is touched, so a conflict or a
missing target leaves the whole tree where it was: still on the source
branch, nothing checked out and nothing merged. A fast-forward is reported
as one, which is usually why a repository looks untouched afterwards.
\texttt{--into} never creates its target: a branch that is not there is as
likely to be a typing mistake as a new branch, so it is named and the
command stops. It cannot be combined with \texttt{--resolve}, which promises
the opposite about a conflict.

\texttt{checkout} warns when it is about to install a different build of the
tool, and still performs the checkout --- looking at an older branch is
legitimate, being surprised by it is not.

A \texttt{--private} or \texttt{--all} \texttt{--into} call only ever
touches part of the tree, and the rest is meant to still be somewhere else
until it too is asked to move. So a second scoped call --- \texttt{--private}
after the plain form, or \texttt{--all} after either --- finishes what the
first one left behind, rather than refusing those repositories for not
having moved yet.

\texttt{--private} merges the writable configuration repositories instead
of the project's own. Read-only ones are never merged.

\begin{center}
\emph{A conflict in any repository leaves no repository merged.}
\end{center}

Every repository in scope is checked before any is merged, using a
read-only test that touches neither worktree nor index. A tree-wide merge
that stopped halfway would leave a workspace no \texttt{.gts} describes
and no command undoes, so it is prevented rather than reported.
\texttt{--dry-run} prints the plan --- and, importantly, the branch each
repository would actually merge, and whether it would merge at all.

\paragraph{The refusal names the files.} When \cgscli{} refuses, it names
every blocked repository and, under each, every file that conflicts:

\begin{lstlisting}[language=bash]
merge refused; no repository was merged: ComplexGitSync: tests/unit/test_documents.py
\end{lstlisting}

A file that both branches changed is only listed when it genuinely
conflicts. Binary files count: a tracked PDF changed on both sides
conflicts just as a source file does, and is reported the same way.
\texttt{--dry-run} shows the same list without merging.

\paragraph{\texttt{--resolve}: fix the conflict instead of reading about
it.} \texttt{--resolve} merges one repository at a time and stops at the
first that conflicts, then opens that repository in your merge tool.

\begin{center}
\emph{\texttt{--resolve} gives up the guarantee above.}
\end{center}

Repositories merged before the conflict stay merged, so the tree can be
left partly merged. That is the price of reaching a conflicted worktree at
all, and it is why this is opt-in: the command prints the warning before it
writes anything, then names what it merged, where it stopped, and which
repositories it never reached.

Your own \texttt{merge.tool} is used whenever you have configured one.
Otherwise VS Code is suggested when it is available, passed for that one
call --- your Git configuration is never written. With no merge tool
available at all, the command prints the command to run by hand and stops,
rather than failing.

\paragraph{Check out the target first.} \texttt{merge} brings a branch
\emph{into} the one the tree is on, exactly as \texttt{git merge} does.
Running \texttt{cgitsync merge multi-branch} while the tree is still on
\texttt{multi-branch} would merge it into itself --- which Git reports as
success while doing nothing --- so \cgscli{} refuses and says to check out
the target branch first.

Landing a feature branch is therefore two commands:

\begin{lstlisting}[language=bash]
$ pixi run cgitsync merge multi-branch
$ pixi run cgitsync merge --private multi-branch
\end{lstlisting}

\subsection{\texttt{pull --private}}

\begin{lstlisting}[language=bash]
$ pixi run cgitsync pull --private
\end{lstlisting}

The other direction. While a project feature branch is open, work lands on
a private/local repository's base branch and its derived branch does not
see it. This fetches and merges the base branch into each one, so a
settings branch does not drift behind the project's. A repository already
on its base branch has nothing to take.

% -----------------------------------------------------------------------
\section{Document Formats}

\texttt{.cgs} means \textbf{ComplexGitSync} specification, \texttt{.gts} means
\textbf{GitTreeState} snapshot, and \texttt{.lgr} means
\textbf{Local Git Register}.

\subsection{Local Authoring Spec (\texttt{.cgs})}

A TOML file that describes the repository tree to manage. It is the only
hand-authored input; it is never modified at runtime.  The normal form has two
top-level keys:

\begin{lstlisting}
project = "CGSil1"
repos = [
  "gitlab:CGS_test/CGSil1",
  "gitlab:CGS_test/CGSil2",
  "github:flipoyo/CGSih1",
]
\end{lstlisting}

Repository identifiers have the form
\texttt{provider:owner/repository}. Nested GitLab namespaces are supported;
the final path segment is the repository name. The document pipeline is:

\begin{center}
\texttt{PARSE -> NORMALIZE -> VALIDATE -> canonical CgsDocument}
\end{center}

The canonical providers and hosts are:

\begin{center}
\begin{tabular}{lll}
\textbf{Provider} & \textbf{Host} & \textbf{Remote generation} \\
\hline
\texttt{github}   & \texttt{github.com}   & SSH and HTTPS \\
\texttt{gitlab}   & \texttt{gitlab.com}   & SSH and HTTPS \\
\texttt{codeberg} & \texttt{codeberg.org} & SSH and HTTPS \\
\texttt{custom}   & explicit URL required & SSH and HTTPS \\
\end{tabular}
\end{center}

Repository-ID parsing is provider-independent and deterministic. For example:

\begin{lstlisting}
codeberg:GX4G/GX4G
\end{lstlisting}

normalizes to \texttt{gitprovider = codeberg},
\texttt{project\_owner\_name = GX4G}, and
\texttt{repo\_name = GX4G}. A custom identifier must use an inline table so
the host is explicit:

\begin{lstlisting}
{ repository = "custom:team/tool", gitprovider_url = "https://git.example.com" }
\end{lstlisting}

No host is inferred for \texttt{custom}. The configured access protocol then
selects \texttt{git@host:owner/repository.git} or
\texttt{https://host/owner/repository.git}.

The textual grammar and its only parser, \texttt{parse\_repo\_id()}, live in
\texttt{cgs\_format.py}. The provider model in \texttt{git\_repo.py} receives the
already-separated provider, owner, and repository fields and composes remotes;
it does not parse authoring shorthand.

The \texttt{.cgs} format pipeline is deterministic and offline. A syntactically
valid repository may therefore reference a host, owner, branch, or tag that is
not currently reachable. Remote existence and reference availability are checked
only by explicit Git operations during runtime resolution and execution.

Normalization supplies \texttt{default\_branch = main} (see ``Branches in a
\texttt{.cgs}'' below for the full chain), makes
\texttt{fallback\_branch} equal to that default, selects
\texttt{access\_protocol = ssh}, enables \texttt{nested\_config = auto}, and
uses the repository name as \texttt{relative\_path}. If exactly one repository
name matches \texttt{project}, its path is inferred as \texttt{.}.

Advanced inline tables are used only where a value differs from those
defaults:

\begin{lstlisting}
project = "demo"
repos = [
  "github:owner/demo",
  { repository = "github:owner/docs", relative_path = "documentation", access_protocol = "https" },
]
\end{lstlisting}

The legacy advanced \texttt{[project]} and \texttt{[[repos]]} tables remain
accepted for project-wide branch settings, custom provider URLs, tags, and
other explicit overrides.  Unknown providers, malformed identifiers,
duplicates, and missing \texttt{project} or \texttt{repos} are rejected before
a \texttt{GitTree} is built.

Two checked-in files document this boundary. Copy
\texttt{examples/template.cgs} for normal authoring. Its equivalent
\texttt{examples/normalized\_template.cgs} expands every inferred field and is
provided as a developer reference for the canonical \texttt{CgsDocument}
shape. The two files are regression-tested for semantic equivalence.

The repository also keeps three checked-in, didactically ordered tutorials
under \texttt{tutorials/} --- easiest first, see
\texttt{tutorials/README.md} for the guided order --- exercising this
authoring surface end to end:
\texttt{tutorials/01\_first\_multi\_repo\_workspace.md} walks a fully
exercised local sandbox topology (\texttt{CGSil1}) through the complete CLI
lifecycle;
\texttt{tutorials/02\_onboarding\_a\_real\_build\_tree.md} hand-authors
a \texttt{.cgs} for the larger, real-world \texttt{cawaqs} tree and the
point where \cgscli{} hands off to the project's existing
\texttt{make}-based build workflow; and
\texttt{tutorials/03\_adopting\_a\_real\_project.md} covers the
real-world \texttt{cawaqsviz} project (not synthetic), which has no
\texttt{.cgs} of its own and still tracks three repositories on two levels
as git submodules,
as one verified procedure: clone, check out the submodules, adopt the whole
tree with \texttt{init-from-submodules} (\texttt{discover} +
\texttt{initialise} + \texttt{import-submodules}, in that order), then
\texttt{branch}/\texttt{checkout}/\texttt{add}/\texttt{commit}/\texttt{push}.

\texttt{examples/cawaqsviz.cgs}, the hand-authored alternative mentioned in
\texttt{tutorials/03\_adopting\_a\_real\_project.md} above,
describes \url{https://gitlab.com/cawaqs/gviz/cawaqsviz} and exercises
several patterns above at once: a nested GitLab subgroup path
(\texttt{cawaqs/gviz/cawaqsviz}, three segments rather than a plain
\texttt{owner/repository} pair), an explicit \texttt{relative\_path = "."}
on the root entry rather than relying on the project-name auto-mount rule,
two GitHub children mounted at non-default relative paths
(\texttt{external/HydrologicalTwinAlphaSeries} and
\texttt{docs/CWV\_user\_guide}). Neither child has a \texttt{.cgs} of its
own, but no override is needed for that: the default \texttt{auto} nested
discovery resolves a repository with zero \texttt{*.cgs} matches as a
normal leaf.
This file previously shipped with a nonexistent repository identifier and
an invalid \texttt{default\_branch}; it was corrected and given real
end-to-end test coverage
(\texttt{tests/integration/test\_cgsi\_topology.py}, and one verified run
against the live repositories).
See \texttt{tutorials/03\_adopting\_a\_real\_project.md}
for the full worked example, and \texttt{bootstrap} above for running
\cgspkg{} against it without cloning \cgspkg{} inside the tree it manages.

Generated specifications use the same concise representation. The
\texttt{GitTree.to\_cgs()} convenience method delegates model projection to
\texttt{cgs\_format.py}; neither \texttt{GitTree} nor \texttt{GitRepo} formats
TOML or implements the authoring grammar. Serialization omits every default
that normalization can deterministically reconstruct and retains inline tables
only for exceptional settings such as a non-default branch, path, protocol,
tag, custom host, or discovery rule.

The supported round trip is semantic:

\begin{center}
\texttt{.cgs -> CgsDocument -> GitTree -> CgsDocument -> .cgs}
\end{center}

Parsing and normalizing the generated file must reproduce the initial
canonical document. Whitespace, comments, and TOML table layout need not be
byte-for-byte identical.

\subsection{Branches in a \texttt{.cgs}}

Four fields decide which branch a repository lands on. Three of them fall
back to each other, so a \texttt{.cgs} that names no branch at all still
sends every repository somewhere.

\begin{center}
\begin{tabular}{lll}
\textbf{Field} & \textbf{Where it is written} & \textbf{Default} \\
\hline
\texttt{project.default\_branch}  & the \texttt{project} table & \texttt{main} \\
\texttt{repos[].default\_branch}  & one repository entry & the project's \\
\texttt{repos[].fallback\_branch} & one repository entry & that repository's own \\
\texttt{repos[].private}           & one repository entry & \texttt{false} \\
\texttt{repos[].writable}         & one repository entry & \texttt{false} \\
\end{tabular}
\end{center}

Read the first three as one chain, longest reach first:

\begin{center}
\texttt{fallback\_branch -> default\_branch -> project.default\_branch -> main}
\end{center}

\texttt{project.default\_branch} is the tree's branch: every repository
inherits it unless the repository says otherwise.
\texttt{repos[].default\_branch} is the branch that one repository targets.
\texttt{repos[].fallback\_branch} is used when the target branch does not
exist on the remote --- \cgscli{} clones the fallback instead and records
that it did so, rather than failing.

Each link is silent. Nothing warns you that a decision was made, so
\textbf{state \texttt{project.default\_branch} in every \texttt{.cgs} you
write}, even when the answer is \texttt{main}. A reader of the file can then
see the branch the tree lands on without knowing what \cgscli{}'s built-in
default happens to be. Every \texttt{.cgs} shipped with ComplexGitSync does
this, and a test enforces it.

\subsubsection{What \texttt{private} means}

A repository is \texttt{private} when it is \textbf{shared with other
projects} and must stay on its own branch.

\begin{center}
\emph{Branch moves stop at a private repository. Tags do not.}
\end{center}

\texttt{cgitsync branch <name>} and \texttt{cgitsync checkout <name>} skip
a private repository and leave it on its declared
\texttt{default\_branch}, because moving it would drag every other project
that mounts it onto your branch.

\subsubsection{Read-only and writable configuration repositories}

A private repository is a \textbf{configuration repository}: shared with
other projects rather than owned by this one. \texttt{private = true} alone
means \textbf{read-only} --- \cgscli{} will not write to it. Add
\texttt{writable = true} for the case where the repository lives outside
your project but the part you use is yours, normally because it sits on a
branch named after your project that nobody else reads:

\begin{lstlisting}
{ repository = "github:you/.sharedSpec", default_branch = "main", private = true },
{ repository = "github:you/.myNotes", default_branch = "MyProject", private = true, writable = true },
\end{lstlisting}

Read-only is the default because the expensive mistake is writing to
something shared by accident, never the reverse. \texttt{writable = true}
on an unprivate repository is rejected at validation: an unprivate repository
is this project's own and is always writable.

This is a \textbf{safety rail, not a permission system}. It stops a
tree-wide sweep from writing where you did not mean; it cannot stop anyone
running \texttt{git} in that directory by hand, and it is no substitute
for branch protection on the remote.

\paragraph{A private/local repository gets a branch per project branch.}
A private/distant repository is private to its owner: you read it, only that
owner writes to it, and nothing you do moves it --- so a branch move leaves
it exactly where it is. A private/local repository is the opposite case:
it holds settings that configure \emph{your} project and you commit to it,
so it needs somewhere to record them per project branch.

\begin{center}
\texttt{<project name>} on \texttt{main}, \quad
\texttt{<project name>\_<branch>} on any other branch
\end{center}

The base is the **project's** name. \texttt{main} takes no suffix, because
the project's main line's settings branch is simply the project's name ---
which is what every existing tree already has, so nothing has to move. The
separator is an underscore because hyphens already occur inside branch
names.

A user never types the derived name. \texttt{cgitsync checkout <branch>}
puts the project's own repositories on that branch and each private/local
one on the branch this rule names, creating it when it is missing.

Without this, one settings branch serves every branch of the project, and
documentation written for an unfinished feature is live on the default
branch describing code that is not there.

\paragraph{\texttt{--private} everywhere.} Every command that touches Git
takes it: \texttt{clone} aside, that is \texttt{pull},
\texttt{pull-force}, \texttt{checkout}, \texttt{branch},
\texttt{close-branch}, \texttt{add}, \texttt{rm}, \texttt{commit},
\texttt{push}, \texttt{merge}, \texttt{tag} and \texttt{freeze}. It
narrows the command to the writable
configuration repositories alone, so they can be worked on without
reaching for plain \texttt{git}. Read-only configuration repositories are
never written to, with or without it. Four of those commands also take
\texttt{--all}, described next.

\paragraph{\texttt{--all}: both halves, one command.} Running every write
twice --- once plainly, once with \texttt{--private} --- is the right
default when the two changes are different changes. When they are one
change, \texttt{--all} does both halves in a single pass, under a single
commit message. Four commands take it: \texttt{add}, \texttt{commit},
\texttt{push} and \texttt{merge}.

\begin{center}
\begin{tabular}{ll}
\textbf{Form} & \textbf{Reaches} \\
\hline
\texttt{cgitsync <cmd>}           & this project's own repositories \\
\texttt{cgitsync <cmd> --private} & the writable configuration repositories \\
\texttt{cgitsync <cmd> --all}     & both of the above, in one pass \\
\end{tabular}
\end{center}

The plain form has not moved, so nothing an existing script runs changes
meaning. \texttt{--all} and \texttt{--private} cannot be combined:
\texttt{--all} already includes everything \texttt{--private} selects. No
form of either writes to a read-only configuration repository --- ``all''
means every repository \cgscli{} may write to, never every repository in
the tree.

The output still names the two halves apart, so giving up the typing does
not mean giving up seeing which repository was written as what. On a tree
with no writable configuration repository, \texttt{--all} does the project
half and reports the other as empty; that is the common case, not a fault.
An explicit \texttt{--private} on the same tree still stops, because there
the user asked for something that is not there.

\begin{center}
\emph{\texttt{merge --all} checks every repository before merging any.}
\end{center}

That is why \texttt{--all} is one pass and never two. Two sequential merges
mean two preflights: the project half would merge, the configuration half
would then refuse, and the tree would be left half-merged --- exactly the
state \texttt{merge}'s whole-scope check exists to prevent.
\texttt{--dry-run} previews the same combined scope, and \texttt{--resolve}
keeps its own warning and its deliberate partial progress.

\paragraph{A pin covers everything inside it.} A configuration repository
often holds repositories of its own, reached through its nested
\texttt{.cgs}. Those are just as shared as the repository holding them, so
they inherit its state: \textbf{a repository inside a read-only
configuration repository is read-only, even when its own entry says
nothing about it}. One inside a writable configuration repository is
writable, and \texttt{--private} reaches it.

A nested entry may restrict itself further --- writing \texttt{private =
true} on its own line, with no \texttt{writable}, makes it read-only
inside a writable parent. It cannot go the other way: \texttt{writable =
true} inside a read-only configuration repository is ignored, because a
repository can never be more open than the one that contains it.

\subsubsection{Which repositories each command touches}

\begin{center}
\begin{tabular}{lll}
\textbf{Command} & \textbf{Reaches} & \textbf{Why} \\
\hline
\texttt{add}, \texttt{commit}, \texttt{push}    & this project's own     & they write this project's history \\
the same, \texttt{--private}                   & writable config repos  & shared, but on your own branch \\
\texttt{branch}, \texttt{checkout}              & this project's own     & a branch must not spread \\
\texttt{tag}, \texttt{freeze-release}           & everything writable    & they create \emph{and push} a tag \\
\texttt{clone}, \texttt{pull}, \texttt{status}  & everything             & reading a shared repo is the point \\
\end{tabular}
\end{center}

\paragraph{Preflight follows the same scope.} Each command checks only
the repositories it is about to touch. A configuration repository that is
behind its upstream, or sitting on its own branch, never blocks a
\texttt{commit} of this project's own work --- that repository is not part
of the commit. Under \texttt{--private} it is in scope, and its branch is
then checked against the branch it is \emph{private} to, not against the
one the project is on.

\texttt{--private} is exclusive, not additive: it selects \emph{only} the
writable configuration repositories, so a shared repository gets its own
command and its own commit message rather than being swept into this
project's. Asking for it in a tree with no writable configuration
repository is an error naming the read-only ones, not a silent no-op. Every
scoped command reports what it left alone:

\begin{lstlisting}[language=bash]
scope=project skipped=5 configuration repo(s) (.claude, .localSpec with --private)
\end{lstlisting}

\texttt{tag} and \texttt{freeze-release} stop at read-only repositories
because they push. Nothing is lost: a \texttt{.gts} snapshot records every
repository's exact \texttt{commit\_sha}, read-only ones included, so an
exact tree is rebuilt from the snapshot rather than from tags.

In practice shared config mounts are usually dot-named
(\texttt{.agentSpec}, \texttt{.localSpec}), and \texttt{cgitsync discover}
drafts a dot-named repository with \texttt{private = true} for that reason.
The rule is \emph{shared}, not \emph{hidden}, though: a plainly named
repository that two projects mount should be private too.

\textbf{The branch a repository is private to decides who sees your commits.}
A mount private to a branch named after your own project is yours to push
freely. A mount private to \texttt{main} is read by every project that mounts
it, so a commit there is published the moment it is pushed --- no branch, no
pull request, no review. Check which of the two you are looking at before
committing to a private mount.

\subsubsection{The two meanings of ``force''}

\texttt{pull-force} is destructive: it discards local changes.
\texttt{--force-protocol} is not: it only switches a remote between
\texttt{ssh} and \texttt{https}. The two are unrelated despite the shared
word.

\subsection{Repository Placement and Nested Discovery}

For every non-root entry, \texttt{relative\_path} is relative to the parent
repository represented by the current \texttt{.cgs}; it is not relative to
the process working directory. At the top level that parent is the project
root. While expanding a nested config, it is the repository containing that
config. The default is the repository name. A unique repository matching the
project name uses \texttt{.} because it represents the current root.

The \texttt{nested\_config} option controls whether traversal continues after
the repository is available:

\begin{itemize}
  \item \texttt{auto} (default) loads the sole root-level \texttt{*.cgs} file,
        if one exists; finding none resolves the repository as a normal
        leaf; multiple matches are rejected as ambiguous.

        A consequence worth knowing before you publish a repository: \textbf{a
        repository that other people may mount inside their tree must keep
        exactly one \texttt{.cgs} at its root}, because \texttt{auto} has no
        way to choose between two. The tree's own root is exempt --- it is
        never scanned for its own \texttt{.cgs}, since that would recurse ---
        so a root holding two \texttt{.cgs} files works fine until the day
        someone mounts it. The error names both files and how to resolve it.
  \item \texttt{disabled} does not inspect the repository for nested config.
  \item A relative path ending in \texttt{.cgs} loads that exact file inside
        the repository, failing if it does not exist. It may not escape the
        repository root.
\end{itemize}

For the integration topology, \texttt{CGSil2.cgs} contains:

\begin{lstlisting}
{ repository = "github:flipoyo/CGSih1", relative_path = "../CGSih1" }
\end{lstlisting}

Because the config describes the tree below \texttt{CGSil2}, the path resolves
from that directory:

\begin{center}
\texttt{<CGSHOME>/CGSil2/../CGSih1 -> <CGSHOME>/CGSih1}
\end{center}

It therefore references the existing sibling declared by \texttt{CGSil1.cgs},
instead of cloning another \texttt{CGSih1} under \texttt{CGSil2}. Discovery
recognizes the already-registered absolute path and keeps the canonical
entry --- this guard runs before \texttt{auto} would ever get a chance to
reopen \texttt{CGSih1.cgs} through the cross-reference, so no
\texttt{nested\_config} override is needed here.

\subsection{Git Tree State Snapshot (\texttt{.gts})}

A TOML file generated automatically by bootstrapping and release operations.
It must never be hand-edited. It captures the exact commit SHAs and lifecycle
state of the tree at a point in time.

\textbf{Top-level sections:}
\begin{itemize}
  \item \texttt{[document]} — \texttt{format\_version}, \texttt{schema\_version},
        \texttt{generated\_at}, \texttt{command\_origin},
        \texttt{hash\_algorithm}, \texttt{snapshot\_hash}.
  \item \texttt{[project]} — project name and \texttt{root\_absolute\_path}.
  \item \texttt{[tree\_state]} — \texttt{lifecycle\_state}, \texttt{is\_ready},
        \texttt{registry\_complete}.
  \item \texttt{[[repo\_state]]} — one entry per repo; includes
        \texttt{commit\_sha}, \texttt{sync\_state}, and ref information.
\end{itemize}

\textbf{Why both version fields and hash metadata exist:}
\begin{itemize}
  \item \texttt{format\_version} identifies the broad document family and keeps
        long-range compatibility intent explicit.
  \item \texttt{schema\_version} identifies the exact field-level contract used
        by validation and canonical snapshot serialization.
  \item \texttt{snapshot\_hash} (with \texttt{hash\_algorithm = "sha256"})
        provides deterministic workspace identity: identical tree state produces
        identical digest even if volatile metadata (\texttt{generated\_at},
        \texttt{command\_origin}) differs.
\end{itemize}

\textbf{How this is used at runtime:}
\begin{itemize}
  \item During snapshot generation, ComplexGitSync writes
        \texttt{schema\_version}, \texttt{hash\_algorithm}, and canonical
        \texttt{snapshot\_hash}.
  \item During load/validation, \texttt{snapshot\_hash} is recomputed from the
        canonical payload and must match.
  \item The project-local \texttt{.lgr} register uses this canonical hash to
        deduplicate equivalent snapshots.
\end{itemize}

\subsection{Local Git Register and Sync Ledger (\texttt{.lgr})}

A TOML file maintained per project at \texttt{<project-root>/<Project\_name>.lgr}.
It is updated automatically whenever a \texttt{.gts} snapshot is written.

\textbf{Sections:}
\begin{itemize}
  \item \texttt{[register]} — current snapshot pointer (\texttt{current\_snapshot\_id},
        \texttt{current\_snapshot\_hash}, \texttt{current\_snapshot\_path}).
  \item \texttt{[[snapshots]]} — list of stable \texttt{gts-XXXXXX} identifiers,
        one per unique workspace state, deduplicated by SHA-256 hash.
  \item \texttt{[[ledger]]} — append-only list of synchronisation events forming
        a DAG via \texttt{parent\_sync\_ids}.
\end{itemize}

\textbf{Ledger event schema:}

\begin{lstlisting}[language=bash]
[[ledger]]
sync_id         = "lgr-000001"
parent_sync_ids = []
operation       = "clone"
timestamp       = "2026-05-20T19:48:50.159Z"
actor           = "username"
workspace_hash  = "5f7c8a2d..."   % 64-character SHA-256 hex digest
gts_snapshot_id = "gts-000001"
affected_repos  = ["demo", "dep-a"]
\end{lstlisting}

The \texttt{workspace\_hash} field equals the canonical SHA-256 digest
(\texttt{document.snapshot\_hash}) of the associated \texttt{.gts} file,
creating a direct link between each ledger event and the immutable workspace
state it records.  Events are parents-before-children in DAG order.

% -----------------------------------------------------------------------
% worked_examples.tex — step-by-step runbooks for the checked-in examples/*
% topologies, ordered from the easiest to the most advanced. Companion to,
% not a repeat of, the per-command reference in Text/user_guide.tex and the
% narrative tutorials under docs/tutorials/: this section gives the exact
% command sequence for each topology and explains what changes between
% them, and points to those other documents for flag semantics and full
% expected-output transcripts.
% -----------------------------------------------------------------------
\section{Worked Examples: From \texttt{CGSil1} to \texttt{cawaqsviz}}
\label{sec:worked-examples}

\cgspkg{} ships four checked-in \texttt{examples/} topologies, and a
same-numbered tutorial under \texttt{docs/tutorials/} for each level below
(\texttt{docs/tutorials/README.md} lists all three in order). Run through
them in order: each one adds exactly one new idea on top of the last, so
together they form a guided path from the simplest possible \texttt{.cgs}
to a real project with no \texttt{.cgs} of its own at all.

\textbf{Every command below is a Pixi task.} \cgscli{} is not a globally
installed executable --- always run \texttt{pixi run cgitsync ...}, never a
bare \texttt{cgitsync ...} (see the \S\ref{chap:user-guide} ``CLI
Overview'' note above).

\begin{center}
\begin{tabular}{lllp{5.3cm}}
\textbf{Level} & \textbf{Topology} & \textbf{Entry point} & \textbf{New idea introduced} \\
\hline
1 & \texttt{CGSil1} & hand-authored \texttt{.cgs} & minimal spec, name-matching auto-mount \\
2 & \texttt{cawaqs} & hand-authored \texttt{.cgs} & the same minimal-field style as Level~1, at 19-repo scale \\
3a & \texttt{cawaqsviz} & hand-authored \texttt{.cgs} & explicit overrides on a real, irregular tree \\
3b & \texttt{cawaqsviz} & \texttt{discover} & drafting a \texttt{.cgs} from a checkout, no authoring at all \\
3c & \texttt{cawaqsviz} & \texttt{import-submodules} & migrating an existing git-submodules project \\
\end{tabular}
\end{center}

Levels~3a--3c are not three different projects: they are three different
\emph{starting points} for the same real project, \texttt{cawaqsviz}, chosen
so this section can demonstrate every route \cgspkg{} offers for arriving
at a working \texttt{.cgs} --- write it by hand, derive it from a checkout,
or migrate it from an existing submodules setup --- without inventing a
fourth toy topology to do it. It is the most advanced level not because
\texttt{cawaqsviz} is a larger tree than \texttt{cawaqs} (it is much
smaller, three repos against nineteen) but because, unlike either earlier
level, it has no \texttt{.cgs} anywhere upstream to copy from.

% -----------------------------------------------------------------------
\subsection{Level 1 --- \texttt{CGSil1}: the minimal synthetic topology}

\texttt{CGSil1} is the easiest way into \cgspkg{}: a 4-repository,
mixed-provider tree with every field left at its default. Its full
\texttt{.cgs} and the rationale for each line are in \S\ref{chap:user-guide}
``Local Authoring Spec'' above;
\texttt{docs/tutorials/01\_first\_multi\_repo\_workspace.md} carries the
full transcript with expected output for every step. What follows is the
command sequence itself.

\begin{lstlisting}[language=bash]
# 1. Fetch the root repo and cgitsync itself (nested mode: cgitsync is
#    cloned inside the tree it will manage).
$ git clone https://gitlab.com/CGS_test/CGSil1
$ cd CGSil1 && git clone https://github.com/flipoyo/ComplexGitSync
$ cd ComplexGitSync

# 2. Check the spec offline, no network beyond parsing.
$ pixi run cgitsync validate ../CGSil1.cgs

# 3. Preview the topology before touching disk.
$ pixi run cgitsync view-tree ../CGSil1.cgs

# 4. Clone every declared repo and reach READY.
$ pixi run cgitsync initialise ../CGSil1.cgs

# 5. The ordinary git cycle, run tree-wide from here on --- no path
#    argument needed, the .gts snapshot is discovered automatically.
$ pixi run cgitsync pull
$ pixi run cgitsync add
$ pixi run cgitsync commit "my commit message"
$ pixi run cgitsync push

# 6. Minimalist release: stage, commit, pull, push, freeze in one call.
$ pixi run cgitsync freeze-release v1.1.0 "release v1.1.0"

# 7. Return to a frozen release at any later point.
$ pixi run cgitsync launch-release v1.1.0
\end{lstlisting}

The only \texttt{.cgs} idea this level relies on is the
\textbf{name-matching auto-mount}: because \texttt{CGSil1}'s repository
identifier's last segment equals \texttt{project = "CGSil1"}, the root is
inferred at \texttt{relative\_path = "."} with no override written anywhere.
That convenience is exactly what stops being safe once a project's
identifier and display name diverge --- which is where Level~3a picks up.

% -----------------------------------------------------------------------
\subsection{Level 2 --- \texttt{cawaqs}: the same minimal style, at scale}

\texttt{cawaqs} scales Level~1's minimal, no-override authoring style up to
a real 19-repository build tree (root + 17 physics libraries across five
GitLab groups + one required build-tooling repo), instead of a synthetic
4-repo one. Every entry's on-disk layout already matches the default
\texttt{relative\_path = <repo\_name>}, so not a single \texttt{relative\_path}
override is written anywhere in \texttt{examples/cawaqs.cgs} --- only the
shared branch-fallback convention. Full narrative in
\texttt{docs/tutorials/02\_onboarding\_a\_real\_build\_tree.md}.

\begin{lstlisting}[language=bash]
$ pixi run cgitsync validate examples/cawaqs.cgs
$ pixi run cgitsync bootstrap examples/cawaqs.cgs cawaqs-smoke-test
$ pixi run cgitsync view-tree --search-dir <path bootstrap printed>

# Move the whole 19-repo tree to a shared branch in one call, falling
# back to main per-repo wherever that branch does not exist.
$ pixi run cgitsync checkout my-feature-branch
\end{lstlisting}

Reaching \texttt{READY} here only replaces the \emph{repository-fetching and
branch-selection} half of \texttt{cawaqs}'s own
\texttt{make\_Cawaqs.sh}/\texttt{make\_Cawaqs\_from\_branches.sh} scripts.
\cgspkg{} does not compile anything: the remaining
\texttt{make -f Makefile} build step is still the user's job, and only
succeeds from this tree with \texttt{LIB\_HYDROSYSTEM\_PATH} left unset or
pointed at the bootstrapped root --- the alternative shared, decoupled
install layout those scripts also support is out of scope for \cgspkg{}'s
containment model.

\texttt{discover} (Level~3b below) is not a route into \texttt{cawaqs.cgs}:
the 17 library repositories never coexist inside one directory tree until
\emph{after} a \texttt{.cgs} already lists them, so there is no single
checkout for a filesystem walk to scan. Authoring this one from documented
developer knowledge, the way Level~3a does for \texttt{cawaqsviz}, is the
only available route.

% -----------------------------------------------------------------------
\subsection{Level 3 --- \texttt{cawaqsviz}: three ways to the same \texttt{.cgs}}

\texttt{cawaqsviz} (\url{https://gitlab.com/cawaqs/gviz/cawaqsviz}) is a
real GitLab project, nested three path segments deep, with two GitHub
children that were, until recently, plain git submodules. It has no
\texttt{.cgs} of its own anywhere upstream. That makes it a good vehicle for
showing all three ways \cgspkg{} can produce a working \texttt{.cgs} for a
project it does not already control the source of, ordered here by how much
the operator has to already know about the tree before running the command.
Full narrative, as one combined ten-step procedure rather than three
independent routes, in \texttt{docs/tutorials/03\_adopting\_a\_real\_project.md}.

\subsubsection{3a --- Write it by hand}

The most explicit, and most error-prone, route: read the project's own
\texttt{.gitmodules} and topology, then author \texttt{examples/cawaqsviz.cgs}
directly. Its corrected content and the exact mistakes the first draft made
(wrong nested-subgroup identifier, missing \texttt{relative\_path}
overrides) are already spelled out in \S\ref{chap:user-guide} ``Local
Authoring Spec''; this level is about running it, not re-deriving it:

\begin{lstlisting}[language=bash]
$ pixi run cgitsync validate examples/cawaqsviz.cgs
$ pixi run cgitsync bootstrap examples/cawaqsviz.cgs cawaqsviz-demo
$ pixi run cgitsync view-tree --search-dir <path bootstrap printed>
\end{lstlisting}

\texttt{bootstrap} is used instead of \texttt{initialise} here because, unlike
\texttt{CGSil1} above, this example is not meant to have \cgspkg{} cloned
inside the tree it manages --- see \texttt{bootstrap} in
\S\ref{chap:user-guide} for why the two commands pick different default
locations.

\subsubsection{3b --- Derive it with \texttt{discover} (autodiscover mode)}

The opposite route: start from nothing but a checkout, and let \cgspkg{}
read the filesystem instead of writing the \texttt{.cgs} by hand. This is
the route to prefer whenever a checkout is already available, since it
cannot repeat the hand-authoring mistakes of 3a --- \texttt{relative\_path}
falls out of the walk directly instead of being typed.

\begin{lstlisting}[language=bash]
# A plain clone leaves submodule paths empty; initialise them first so
# discover has something to find at those paths. --init only (not
# --recursive): a nested submodule one level deeper would otherwise show
# up as a fourth repository discover has no use for here. (Want that
# nested submodule too? Use --init --recursive here and
# import-submodules --recursive --apply in 3c.)
$ git clone https://gitlab.com/cawaqs/gviz/cawaqsviz cawaqsviz
$ cd cawaqsviz && git submodule update --init

# Dry run: report what discover sees, write nothing.
$ pixi run cgitsync discover . --max-depth 3

# Satisfied with the report: draft the .cgs.
$ pixi run cgitsync discover . --write cawaqsviz-discovered.cgs
$ pixi run cgitsync validate cawaqsviz-discovered.cgs
\end{lstlisting}

The draft reconstructs 3a's hand-corrected file field-for-field: the root at
\texttt{relative\_path = "."} and both children at their real submodule
paths (\texttt{external/HydrologicalTwinAlphaSeries},
\texttt{docs/CWV\_user\_guide}) rather than their bare repo names. Neither
child has a \texttt{.cgs} of its own, but \texttt{discover} writes no
\texttt{nested\_config} override for that: the default \texttt{"auto"}
already resolves a repository with zero nested \texttt{*.cgs} matches as a
normal leaf. See \texttt{discover} in \S\ref{chap:user-guide} for exactly
what is and is not reported (no network calls, unresolvable remotes are
warned about and left out, never guessed at).

\subsubsection{3c --- Migrate it with \texttt{import-submodules}}

The historical route: before either \texttt{.cgs} above existed,
\texttt{cawaqsviz} tracked both children as real git submodules. This mode
starts from that state and converts it, rather than describing a tree that
is assumed to already be plain clones.

\begin{lstlisting}[language=bash]
# Dry run against a real cawaqsviz checkout that still has submodules:
# reports name, path, URL, branch per submodule; changes nothing.
$ pixi run cgitsync import-submodules /path/to/cawaqsviz

# Convert: drop the gitlinks, retire .gitmodules, extend .gitignore.
$ pixi run cgitsync import-submodules /path/to/cawaqsviz --apply

# Review and commit the resulting working-tree changes as usual.
$ cd /path/to/cawaqsviz && git status
$ git commit -m "chore: retire git submodules in favour of ComplexGitSync"
\end{lstlisting}

This is the whole job: turning gitlinks into plain clones on disk. It does
not also write a \texttt{.cgs} --- \texttt{.gitmodules} never records the
root's own identity, and a checkout worth converting already has a
\texttt{.cgs} (3a) or can get one from \texttt{discover} (3b), run before
or after \texttt{--apply}. The full before/after picture (index,
\texttt{.gitmodules}, \texttt{.gitignore}) and the automated test that
proves it against fixture repositories are in
\texttt{docs/tutorials/03\_adopting\_a\_real\_project.md} \S4 --- this
level is the condensed command sequence, not a restatement of that
walkthrough. \texttt{import-submodules} and \texttt{discover} answer
different questions about the same tree: \texttt{discover} reports what is
\emph{checked out}, \texttt{import-submodules} reads what git
\emph{declares} in \texttt{.gitmodules} and acts on it --- a submodule path
that was never initialised is invisible to 3b but not to 3c.

\subsubsection{A shared caveat across all three}

\cgspkg{} stores no credentials and has no private-repository
authentication story (\S\ref{chap:user-guide} ``\texttt{import-submodules}''
and ``\texttt{discover}''). All three routes above assume every repository
in the tree is reachable anonymously or through credentials the ambient
environment (\texttt{ssh-agent}, an HTTPS credential helper) already holds.


% -----------------------------------------------------------------------
\section{Configuration and Environment}

\cgspkg{} uses project files for topology and does not require global topology
configuration. Snapshot discovery can use \texttt{CGSHOME}; when it is not set,
the \texttt{initialise(.cgs)} output-path default is \texttt{CGSPATH=../..}
(nested mode: \cgspkg{} cloned inside the tree it manages); later
commands can also walk upward from the current directory to find
\texttt{.cgitsync/state/}. \texttt{bootstrap(.cgs)} uses a different
default instead --- \texttt{CGSPATH=\$HOME/.cgs/CGS<timestamp>/}, created
automatically --- since it is meant for a \cgspkg{} clone that is not
nested inside any project tree; see \texttt{bootstrap} above. The logging
and runtime-state subsystems use the user state directory by default:

\begin{itemize}
  \item logs: \texttt{~/.local/state/ComplexGitSync/logs/}
  \item latest runtime-state mapping:
        \texttt{~/.local/state/ComplexGitSync/runtime/}
\end{itemize}

If \texttt{XDG\_STATE\_HOME} is set, these directories are rooted there
instead.

% -----------------------------------------------------------------------
\section{Error Reference}

\begin{description}
  \item[\texttt{ConfigValidationError}] Raised when \texttt{project} or
    \texttt{repos} is missing, a repository identifier is malformed, or a
    canonical field is invalid (for example an unknown \texttt{gitprovider}).
  \item[\texttt{GitSyncError}] Raised when clone-time Git operations fail,
    for example because the remote branch cannot be resolved or the target
    directory is already occupied.
  \item[Exit code 1] General CLI failure (parse error, validation error).
  \item[Exit code 2] Command exists but is not yet implemented.
\end{description}

% -----------------------------------------------------------------------
\section{Troubleshooting}

\begin{description}
  \item[``\texttt{gitprovider\_url} is required for custom provider'']
    Add a \texttt{gitprovider\_url} field to the offending
    \texttt{[[repos]]} entry in your \texttt{.cgs} file.
  \item[\texttt{validate} exits non-zero despite a correct-looking file]
    Check that \texttt{project} is present and every \texttt{repos} entry uses
    \texttt{provider:owner/repository}. For advanced \texttt{[[repos]]}
    blocks without a \texttt{repository} shorthand, include
    \texttt{project\_owner\_name} and \texttt{project\_name}.
  \item[\texttt{Clone destination already exists and is not empty}]
    Re-run \texttt{clone} with \texttt{--target-dir} to force a clean location,
    or let \cgscli{} choose a suffixed default directory.
  \item[\texttt{No cloneable branch found for <repo>}]
    Confirm that the remote exposes either the repo's
    \texttt{default\_branch} or its \texttt{fallback\_branch}.
  \item[PyYAML import error when calling \texttt{from\_yaml}/\texttt{to\_yaml}]
    Add PyYAML in the Pixi environment: \texttt{pixi add pyyaml}.
\end{description}
